\documentclass[11pt]{article}

\usepackage[margin=1in]{geometry}
\usepackage{amsmath,amssymb,amsthm,mathtools}
\usepackage{bm}
\usepackage{xcolor}
\usepackage{enumitem}
\usepackage{graphicx}
\usepackage{hyperref}
\usepackage{booktabs}
\usepackage{listings}

\definecolor{leankeyword}{RGB}{32,74,135}
\definecolor{leancomment}{RGB}{78,120,78}
\definecolor{leanstring}{RGB}{150,55,55}
\definecolor{leanbackground}{RGB}{248,248,248}
\lstdefinelanguage{Lean}{
  morekeywords={abbrev,by,case,class,def,deriving,else,end,example,
    export,extends,for,fun,if,import,in,inductive,instance,let,match,
    namespace,noncomputable,open,private,protected,section,set_option,
    simp,structure,then,theorem,variable,where,with},
  sensitive=true,
  morecomment=[l]{--},
  morecomment=[s]{/-}{-/},
  morestring=[b]"
}
\lstset{
  language=Lean,
  basicstyle=\ttfamily\small,
  keywordstyle=\color{leankeyword},
  commentstyle=\color{leancomment},
  stringstyle=\color{leanstring},
  backgroundcolor=\color{leanbackground},
  frame=single,
  rulecolor=\color{black!20},
  breaklines=true,
  breakatwhitespace=true,
  columns=fullflexible,
  keepspaces=true,
  showstringspaces=false,
  tabsize=2,
  literate=
    {⁻¹}{{$^{-1}$}}2
    {×}{{$\times$}}1
    {ˢ}{{$^{s}$}}1
    {•}{{$\mathbin{\cdot}$}}1
    {∈}{{$\in$}}1
    {∧}{{$\land$}}1
    {⊆}{{$\subseteq$}}1
}

\hypersetup{
  colorlinks=true,
  linkcolor=blue,
  citecolor=blue,
  urlcolor=blue
}

\newtheorem{theorem}{Theorem}[section]
\newtheorem{proposition}[theorem]{Proposition}
\newtheorem{lemma}[theorem]{Lemma}
\newtheorem{corollary}[theorem]{Corollary}
\newtheorem{remark}[theorem]{Remark}
\numberwithin{equation}{section}

\newcommand{\E}{\mathbb{E}}
\newcommand{\Pbb}{\mathbb{P}}
\newcommand{\R}{\mathbb{R}}
\newcommand{\N}{\mathbb{N}}
\newcommand{\1}{\mathbf 1}
\newcommand{\sech}{\operatorname{sech}}
\newcommand{\cP}{\mathcal P}
\newcommand{\cK}{\mathcal K}
\newcommand{\cR}{\mathcal R}
\newcommand{\dd}{\,\mathrm d}
\newcommand{\norm}[1]{\left\lVert #1\right\rVert}
\newcommand{\ip}[2]{\left\langle #1,#2\right\rangle}
\newcommand{\Tr}{\operatorname{Tr}}
\newcommand{\diag}{\operatorname{diag}}
\newcommand{\sgn}{\operatorname{sgn}}
\newcommand{\Var}{\operatorname{Var}}
\newcommand{\Cov}{\operatorname{Cov}}
\newcommand{\e}{\mathrm e}
\newcommand{\braket}[1]{\left\langle #1\right\rangle}

\title{A quantitative replica-symmetric bound in strict AT region}
\author{
  Seiichiro Kusuoka\thanks{Department of Mathematics and Mathematical Sciences, Graduate School of Science, Kyoto University. E-mail: \texttt{kusuoka@math.kyoto-u.ac.jp}}
  \and
  Shuta Nakajima\thanks{Department of Mathematics, Keio University. E-mail: \texttt{njima@keio.jp}}
}
\date{}

\begin{document}
\subsection{Model and the strict AT region}

Let $\Sigma_N\coloneqq\{-1,1\}^N$.  For $\beta>0$ and $h>0$, let
\[
 H_N(\sigma)
 \coloneqq
 \frac{\beta}{\sqrt N}\sum_{1\le i<j\le N}g_{ij}\sigma_i\sigma_j
 +h\sum_{i=1}^N\sigma_i ,
\]
where the $g_{ij}$ are independent standard normal random variables.
For replicas $\sigma^1,\sigma^2,\ldots$ sampled independently from the
Gibbs measure for fixed disorder, put
\[
 R_{ab}\coloneqq\frac1N\sum_{i=1}^N\sigma_i^a\sigma_i^b.
\]

Let $q(\beta,h)$ be the unique solution of
\begin{equation}
 q=\E\tanh^2\bigl(h+\beta\sqrt q\,Z\bigr),
 \qquad Z\sim N(0,1).
 \label{eq:q}
\end{equation}
and write $q\coloneqq q(\beta,h)$ when the parameters are fixed.
Existence and uniqueness for $h>0$ are standard; see, for example,
\cite[Proposition~1.3.8]{Talagrand}.  Continuity follows directly from
uniqueness: if $(\beta_n,h_n)\to(\beta,h)$, every convergent
subsequence of $q(\beta_n,h_n)\in[0,1]$ has, by dominated convergence
in \eqref{eq:q}, a limit solving the fixed-point equation at
$(\beta,h)$.  Uniqueness identifies that limit with $q(\beta,h)$, so
the entire sequence converges.
Set
\begin{equation}
 r\coloneqq\E\tanh^4\bigl(h+\beta\sqrt q\,Z\bigr),
 \label{eq:ra}
\end{equation}
and let
\begin{equation}
 \alpha\coloneqq\beta^2(1-2q+r)
 =\beta^2\E\sech^4\bigl(h+\beta\sqrt q\,Z\bigr).
 \label{eq:alpha}
\end{equation}
We work on a fixed compact set
\begin{equation}
 \cK\Subset\{(\beta,h):\beta>0,\ h>0,\ \alpha(\beta,h)<1\}.
 \label{eq:K}
\end{equation}
Thus, for some $\delta_{\cK}>0$,
\begin{equation}
 1-\alpha\ge\delta_{\cK}
 \qquad\text{on }\cK.
 \label{eq:delta}
\end{equation}
The maps $q,r,\alpha$ are therefore continuous on $\cK$.  In
particular, compactness and $h>0$ imply
\begin{equation}
 0<q_{\cK}\coloneqq\inf_{(\beta,h)\in\cK}q
 \le \sup_{(\beta,h)\in\cK}q<1.
 \label{eq:qcompact}
\end{equation}
Indeed, $q=0$ would contradict \eqref{eq:q} and $h>0$, while
$q<1$ because $\tanh^2(h+\beta\sqrt q\,Z)<1$ a.s.

\subsection{The replica-symmetric smart path}

For $0\le s\le1$, define the smart-path Hamiltonian
\begin{equation}
 H_{N,s}(\sigma)
 \coloneqq
 \frac{\beta\sqrt s}{\sqrt N}
 \sum_{1\le i<j\le N}g_{ij}\sigma_i\sigma_j
 +\sum_{i=1}^N
 \bigl(h+\beta\sqrt{(1-s)q}\,z_i\bigr)\sigma_i,
 \label{eq:path}
\end{equation}
where $(z_i)$ is an independent standard Gaussian family.  We write
$\braket{\cdot}_s$ for the Gibbs expectation and
\[
 \nu_s[f]\coloneqq\E\braket{f}_s,\qquad
 Q_{ab}\coloneqq R_{ab}-q.
\]
We also write
\begin{equation}
 \phi_N(s)\coloneqq\frac1N\E\log
 \sum_{\sigma\in\Sigma_N}e^{H_{N,s}(\sigma)}.
 \label{eq:pathfreeenergy}
\end{equation}

\subsection{Main result}

Define
\begin{equation}
 A_s\coloneqq\nu_s[Q_{12}^2],\qquad
 B_s\coloneqq\nu_s[Q_{12}Q_{13}],\qquad
 C_s\coloneqq\nu_s[Q_{12}Q_{34}].
 \label{eq:ABC}
\end{equation}

\begin{theorem}
\label{thm:main}
There is a constant $C_{\cK}<\infty$ such that, for every $N\ge1$,
\[
 \sup_{\substack{(\beta,h)\in\cK\\0\le s\le1}}
 N\,\nu_s[Q_{12}^2]\le C_{\cK}.
\]
Consequently, if
\[
 \phi_N(\beta,h)\coloneqq
 \frac1N\E\log\sum_{\sigma\in\Sigma_N}e^{H_N(\sigma)}
\]
and
\[
 \phi^{\mathrm{RS}}(\beta,h)
 \coloneqq
 \log2+\E\log\cosh(h+\beta\sqrt q\,Z)
 +\frac{\beta^2}{4}(1-q)^2,
\]
then
\[
 0\le\phi^{\mathrm{RS}}(\beta,h)-\phi_N(\beta,h)
 \le\frac{C_{\cK}}N
\]
uniformly on $\cK$.

Moreover, uniformly for $(\beta,h,s)\in\cK\times[0,1]$,
\[
 N\bigl(A_s-2B_s+C_s\bigr)
 =
 \frac{\alpha}{\beta^2(1-s\alpha)}+o_{\cK}(1).
\]
Here and below, $o_{\cK}(1)\to0$ as $N\to\infty$, uniformly over the
 compact set $\mathcal K$.
\end{theorem}

\subsection{Proof strategy}

The proof is divided into the following steps.  First, the scalar
finite-step PDE is solved explicitly, and the strict AT sign is proved
uniformly along the smart path.  Second, the specialized two-replica
Guerra--Talagrand upper bound is established in
Appendix~\ref{app:specialGT}, while the main text derives from it and
the strict AT sign a uniform quadratic bound for the constrained
pressure.  Third, a quadratically
coupled pressure and Gronwall's inequality produce an $o(1)$ overlap
bound, and Gaussian concentration gives fixed-deviation exponential
tails.  Fourth, a cavity expansion gives a three-dimensional linear
system for $A_s,B_s,C_s$ with a cubic remainder.  Uniform invertibility
of the cavity matrix and the fixed-deviation estimate allow that
remainder to be absorbed.  The free-energy estimate follows by
integrating the exact smart-path identity, and the replicon
susceptibility follows from the replicon row of the same cavity system.

\section{Replica-symmetric coercivity}\label{sec:coercivity}

\subsection{Gaussian calculus}\label{subsec:gaussian}

We record Gaussian integration by parts, which is used below.  Let
$G\coloneqq(G_1,\ldots,G_m)$ be centered Gaussian with covariance
matrix $C$.
If $F$ is continuously differentiable with bounded first derivatives,
then Gaussian integration by parts gives
\begin{equation}
 \E[G_iF(G)]=\sum_{j=1}^m C_{ij}\E[\partial_jF(G)].
 \label{eq:GIBP}
\end{equation}

We also use the following standard consequence of Gaussian concentration.
Let \(G\) be a standard Gaussian vector, and suppose that, for each
\(v\in I\), the random variable \(F_v\coloneqq F_v(G)\) is an
\(L\)-Lipschitz
function of \(G\), where \(I\) is finite. Then
\begin{equation}
 \E\max_{v\in I}\{F_v-\E F_v\}
 \le L\sqrt{2\log|I|}.
 \label{eq:gaussianmaxgeneral}
\end{equation}
Indeed, the Gaussian concentration inequality (see, for example,
\cite{Ledoux}) gives, for every \(t>0\)
and every \(v\in I\),
\[
 \E\exp\bigl\{t(F_v-\E F_v)\bigr\}
 \le \exp\left(\frac{t^2L^2}{2}\right).
\]
\[
\begin{aligned}
 \E\max_{v\in I}(F_v-\E F_v)
 &\le\frac1t\log\sum_{v\in I}
 \E e^{t(F_v-\E F_v)}
 \le\frac{\log|I|}{t}+\frac{tL^2}{2}.
\end{aligned}
\]
For $|I|\ge2$ and $L>0$, optimize at
$t=\sqrt{2\log|I|}/L$.  The cases $|I|=1$ and $L=0$ are immediate.


\subsection{The scalar replica-symmetric PDE and the strict AT sign}\label{subsec:ATsign}

For fixed $(\beta,h,s)\in\cK\times[0,1]$, put
\begin{equation}
 x_q(u)\coloneqq\1_{[q,1]}(u),\qquad
 \xi_s(u)\coloneqq\beta^2(1-s)q\,u+\frac{s\beta^2}{2}u^2.
 \label{eq:xqxi}
\end{equation}
Let $\Psi:[0,1]\times\mathbb R\to\mathbb R$ solve
\begin{equation}
 \partial_u\Psi+\frac{s\beta^2}{2}
 \bigl(\partial_{xx}\Psi+x_q(u)(\partial_x\Psi)^2\bigr)=0,
 \qquad \Psi(1,x)=\log\cosh x.
 \label{eq:parisipde}
\end{equation}
This finite-step equation has the explicit semigroup representation
\begin{align}
 \Psi(u,x)
 &=\log\cosh x+\frac{s\beta^2}{2}(1-u), &&q\le u\le1,
 \label{eq:PDEupper}\\
 \Psi(u,x)
 &=\E\Psi\bigl(q,x+\beta\sqrt{s}\sqrt{q-u}\,Z\bigr)\notag\\
 &=\E\log\cosh\bigl(x+\beta\sqrt{s}\sqrt{q-u}\,Z\bigr)
   +\frac{s\beta^2}{2}(1-q), &&0\le u\le q.
 \label{eq:PDElower}
\end{align}
Indeed, on $[q,1]$ the Cole--Hopf transform $\exp\Psi$ solves the
backward heat equation, and on $[0,q]$ the equation itself is the
backward heat equation.  These formulas prove existence, uniqueness,
and all differentiability statements used below.

Let $Z_0$ and a Brownian motion $W$ be independent and define the local
field process by
\begin{equation}
 \dd X_u=s\beta^2x_q(u)\partial_x\Psi(u,X_u)\,\dd u
 +\beta\sqrt{s}\,\dd W_u,
 \qquad X_0=h+\beta\sqrt{(1-s)q}\,Z_0.
 \label{eq:localfield}
\end{equation}
The replica-symmetric trial value along the path is
\begin{equation}
 P_s^*\coloneqq\log2+
 \E\Psi(0,h+\beta\sqrt{(1-s)q}\,Z_0)
 -\frac{s\beta^2}{2}\int_q^1u\,\dd u.
 \label{eq:scalartrial}
\end{equation}
Using \eqref{eq:PDEupper}--\eqref{eq:PDElower} and
$\beta^2(1-s)q+s\beta^2q=\beta^2q$, this reduces to
\begin{equation}
 P_s^*=\log2+\E\log\cosh(h+\beta\sqrt q\,Z)
 +\frac{s\beta^2}{4}(1-q)^2.
 \label{eq:RSpathvalue}
\end{equation}

\begin{lemma}[Diffusion comparison]
\label{lem:diffusioncomparison}
Suppose $s\beta^2(1-q)>1$.  Let $(m_u)_{q\le u\le1}$ solve
\[
 \dd m_u=\beta\sqrt{s}(1-m_u^2)\,\dd W_u,
 \qquad m_q\stackrel d=\tanh(h+\beta\sqrt q\,Z),
\]
with the initial variable independent of the Brownian increments after
time $q$.  Then
\begin{equation}
 \beta^2\E[(1-m_u^2)^2]\le\alpha,qquad q\le u\le1.
 \label{eq:strictATdiffusion}
\end{equation}
\end{lemma}

\begin{proof}
Put $\sigma^2\coloneqq\beta^2q$.  We first show that
$h<\sigma^2$.  The function
$c\mapsto\E\sech^2(c+\sigma Z)$ is nonincreasing on $[0,\infty)$.
Indeed, if $f(x)\coloneqq\sech^2x$ and $\varphi_\sigma$ is the centered
Gaussian density, then for $c\ge0$,
\[
 \frac{\dd}{\dd c}(f*\varphi_\sigma)(c)
 =\int_0^\infty f'(x)
   \{\varphi_\sigma(x-c)-\varphi_\sigma(x+c)\}\,\dd x\le0.
\]
If $h\ge\sigma^2$, the fixed-point equation gives
\[
 1-q=\E\sech^2(h+\sigma Z)
 \le\E\sech^2(\sigma^2+\sigma Z).
\]
For a symmetric $\{-1,1\}$-valued variable $\varepsilon$, set
$Y\coloneqq\sigma\varepsilon+Z$.  Since
$\E[\varepsilon\mid Y]=\tanh(\sigma Y)$, conditional expectation and
the best linear estimator of $\varepsilon$ from $Y$ give
\[
 \E\sech^2(\sigma^2+\sigma Z)
 =\E\bigl[(\varepsilon-\E[\varepsilon\mid Y])^2\bigr]
 \le(1+\sigma^2)^{-1}.
\]
It would follow that $\beta^2(1-q)\le1$, a contradiction.  Hence
$h<\sigma^2$.

Use the time variable $\lambda\coloneqq s\beta^2(u-q)$ and set
$v_\lambda\coloneqq1-m_{q+\lambda/(s\beta^2)}^2$.
Conditional on $m_q=\tanh x$, Girsanov's formula gives
\[
 \E[\psi(X_{q+\lambda/(s\beta^2)})\mid X_q=x]
 =
 \frac{\E[\psi(x+\sqrt\lambda G)\cosh(x+\sqrt\lambda G)]}
      {\E[\cosh(x+\sqrt\lambda G)]},
\]
where $m_u=\tanh X_u$.  Writing $S\coloneqq\sech^2X_q$ and taking
$\psi=\sech^4$ gives
\[
 \E[v_\lambda^2\mid X_q=x]
 =\frac{e^{-\lambda/2}\sech x}{2}\E_G
 \bigl[\sech^3(x+\sqrt\lambda G)
       +\sech^3(x-\sqrt\lambda G)\bigr].
\]
For $y\in\mathbb R$ and $S=\sech^2x$,
\[
 \frac{\sech x}{2}\{\sech^3(x+y)+\sech^3(x-y)\}
 =S^2\cosh y\,
 \frac{1+(4-3S)\sinh^2y}{(1+S\sinh^2y)^3}.
\]
Consequently,
\begin{equation}
 \E v_\lambda^2=\E[S^2F_\lambda(S)],
 \label{eq:Flambda-representation}
\end{equation}
where
\[
 F_\lambda(y)\coloneqq e^{-\lambda/2}\E\left[
 \cosh(\sqrt\lambda G)
 \frac{1+(4-3y)\sinh^2(\sqrt\lambda G)}
      {(1+y\sinh^2(\sqrt\lambda G))^3}\right].
\]
For $t\coloneqq\sinh^2(\sqrt\lambda G)$, put
\[
 H_t(y)\coloneqq\frac{1+(4-3y)t}{(1+yt)^3}.
\]
Then
\[
 H_t'(y)=-\frac{6t\{1+(2-y)t\}}{(1+yt)^4}\le0,
\]
so $F_\lambda$ is decreasing.  Moreover, with $y=1-u^2$,
\[
 \int_0^1H_t(y)\frac{\dd y}{2\sqrt{1-y}}
 =\left[\frac{u}{(1+t-tu^2)^2}\right]_{0}^{1}=1.
\]
Since $e^{-\lambda/2}\E\cosh(\sqrt\lambda G)=1$,
\begin{equation}
 \int_0^1F_\lambda(y)\frac{\dd y}{2\sqrt{1-y}}=1.
 \label{eq:Fmeanone}
\end{equation}

Let $\mu_0(\dd y)\coloneqq\dd y/(2\sqrt{1-y})$ and
$x(y)\coloneqq\operatorname{arcosh}(y^{-1/2})$.  Since
$X_q\sim N(h,\sigma^2)$, the density of
$\E[S^2\1_{\{S\in\dd y\}}]$ with respect to $\mu_0$, up to a positive
constant, is
\[
 \rho(y)\coloneqq y\exp\left\{-\frac{x(y)^2+h^2}{2\sigma^2}\right\}
 \cosh\left(\frac{h x(y)}{\sigma^2}\right).
\]
As a function of $x\ge0$, its logarithmic derivative is
\[
 -2\tanh x-\frac{x}{\sigma^2}
 +\frac{h}{\sigma^2}\tanh\left(\frac{hx}{\sigma^2}\right)<0,
\]
because $h\le\sigma^2$.  Thus $\rho$ is increasing in $y$.
The opposite-monotonicity inequality, followed by
\eqref{eq:Fmeanone}, now gives
\[
 \beta^2\E v_\lambda^2
 \le\left(\int F_\lambda\,\dd\mu_0\right)
       \left(\beta^2\int\rho\,\dd\mu_0\right)
 =\beta^2\E S^2=\alpha.
\]
This is \eqref{eq:strictATdiffusion}.
\end{proof}


\begin{proposition}
\label{prop:pathRS}
Set
\[
 g_s(u)\coloneqq\E\bigl[\partial_x\Psi(u,X_u)^2\bigr].
\]
Uniformly on $\cK\times[0,1]$,
\begin{equation}
 g_s(u)-u
 \begin{cases}
 >0,&0\le u<q,\\
 =0,&u=q,\\
 <0,&q<u\le1.
 \end{cases}
 \label{eq:ATsign}
\end{equation}
Moreover, there are $c_{\cK},\varepsilon_{\cK}>0$ such that
\begin{equation}
 |g_s(u)-u|\ge c_{\cK}|u-q|
 \qquad\text{whenever }|u-q|\le\varepsilon_{\cK}.
 \label{eq:linearATsign}
\end{equation}
The sign is uniform away from $q$ on the compact family.
\end{proposition}

\begin{proof}
Differentiating \eqref{eq:PDEupper} in $u$ and $x$ shows that, for
$u\ge q$,
\begin{align*}
 \partial_x\Psi(u,x)&=\tanh x,
 &\partial_{xx}\Psi(u,x)&=\sech^2x,\\
 \partial_{ux}\Psi(u,x)&=0,
 &\partial_{xxx}\Psi(u,x)&=-2\tanh x\,\sech^2x.
\end{align*}
For $u\le q$, we may differentiate \eqref{eq:PDElower} under the
expectation, since the first two derivatives of $\log\cosh$ are bounded.
Using the preceding identities at time $q$ gives
\begin{align*}
 \partial_x\Psi(u,x)
 &=\E\tanh\bigl(x+\beta\sqrt{s}\sqrt{q-u}\,Z\bigr),\\
 \partial_{xx}\Psi(u,x)
 &=\E\sech^2\bigl(x+\beta\sqrt{s}\sqrt{q-u}\,Z\bigr),\\
 \partial_{ux}\Psi(u,x)
 &=s\beta^2\E\bigl[\tanh(Y_{u,x})\sech^2(Y_{u,x})\bigr],\\
 \partial_{xxx}\Psi(u,x)
 &=-2\E\bigl[\tanh(Y_{u,x})\sech^2(Y_{u,x})\bigr],
 \qquad
 Y_{u,x}\coloneqq x+\beta\sqrt{s}\sqrt{q-u}\,Z.
 \end{align*}
On $[0,q]$ the drift in \eqref{eq:localfield} vanishes, because
$x_q=0$ there up to the immaterial endpoint $q$.  Consequently,
\[
 X_q=X_u+\beta\sqrt{s}(W_q-W_u).
\]
With $\mathcal F_u\coloneqq\sigma(Z_0,W_r:0\le r\le u)$, the increment
$W_q-W_u$ is independent of $\mathcal F_u$ and has law
$\sqrt{q-u}\,Z$.  Substituting $x=X_u$ in the last two semigroup
identities therefore proves
\[
 \partial_x\Psi(u,X_u)=\E[\tanh(X_q)\mid\mathcal F_u],
 \qquad
 \partial_{xx}\Psi(u,X_u)=\E[\sech^2(X_q)\mid\mathcal F_u].
\]
Indeed, on $[0,q]$ the drift-free representation gives
\[
 X_q=h+\beta\sqrt{(1-s)q}\,Z_0+\beta\sqrt{s}\,W_q.
\]
The two Gaussian terms on the right are independent, and their variances
sum to $\beta^2(1-s)q+s\beta^2q=\beta^2q$.  Hence
\begin{equation}
 X_q\stackrel d=h+\beta\sqrt q\,Z,
 \label{eq:Xq}
\end{equation}
as claimed.  To compute $g_s'$, set
$M_u\coloneqq\partial_x\Psi(u,X_u)$.  Differentiating the heat equation on
$[0,q]$ in $x$ and applying Ito's formula yield
\begin{align*}
 \dd M_u
 &=\left(\partial_{ux}\Psi(u,X_u)
   +\frac{s\beta^2}{2}\partial_{xxx}\Psi(u,X_u)\right)\dd u
   +\beta\sqrt{s}\,\partial_{xx}\Psi(u,X_u)\,\dd W_u\\
 &=\beta\sqrt{s}\,\partial_{xx}\Psi(u,X_u)\,\dd W_u.
\end{align*}
A second application of Ito's formula, now to $M_u^2$, gives
\[
 \dd(M_u^2)
 =2\beta\sqrt{s}\,M_u\partial_{xx}\Psi(u,X_u)\,\dd W_u
  +s\beta^2\partial_{xx}\Psi(u,X_u)^2\,\dd u.
\]
All coefficients here are bounded, so the stochastic integral has mean
zero.  Consequently, for $u<q$,
\[
 g_s'(u)
 =s\beta^2\E[\partial_{xx}\Psi(u,X_u)^2].
\]
Moreover, the conditional representation above and conditional Jensen
give
\begin{align*}
 \E[\partial_{xx}\Psi(u,X_u)^2]
 &=\E\left[\E[\sech^2(X_q)\mid\mathcal F_u]^2\right]\\
 &\le \E\left[\E[\sech^4(X_q)\mid\mathcal F_u]\right]
 =\E\sech^4(X_q).
\end{align*}
Together with \eqref{eq:Xq} and \eqref{eq:ra}--\eqref{eq:alpha}, this
proves
\[
 g_s'(u)\le s\beta^2\E\sech^4(X_q)=s\alpha,
 \qquad 0\le u<q.
\]
Since $g_s(q)=q$,
\begin{equation}
 g_s(u)-u\ge(1-s\alpha)(q-u)
 \qquad(0\le u<q).
 \label{eq:leftstrict}
\end{equation}

For $u\ge q$, put $m_u\coloneqq\tanh X_u$.  The drift cancels:
\begin{equation}
 \dd m_u=\beta\sqrt{s}(1-m_u^2)\,\dd W_u.
 \label{eq:mdiffusion}
\end{equation}
Thus, if $f(u)\coloneqq\E m_u^2-u$,
\begin{equation}
 f(q)=0,\qquad f'(u)=s\beta^2\E[(1-m_u^2)^2]-1,
 \qquad f'(q)=s\alpha-1.
 \label{eq:fprime}
\end{equation}
If $s\beta^2(1-q)\le1$, then
$f'(q)=s\alpha-1\le-\delta_{\cK}<0$, so $f<0$ immediately to the
right of $q$.  If $f$ ever returned to zero, let $u_0>q$ be its first
zero after this initial negative interval.  Then differentiability
would give $f'(u_0)\ge0$.  On the other hand,
\[
 \E[(1-m_{u_0}^2)^2]\le1-\E m_{u_0}^2=1-u_0,
\]
so $f'(u_0)\le s\beta^2(1-u_0)-1
<s\beta^2(1-q)-1\le0$, a contradiction.  Hence $f(u)<0$ for every
$u>q$.

If $s\beta^2(1-q)>1$, Lemma~\ref{lem:diffusioncomparison} and
\eqref{eq:fprime} give $f'(u)\le s\alpha-1<0$.  This proves the
right-hand sign in the remaining case.

Finally, \eqref{eq:leftstrict} and
$f'(q)=-(1-s\alpha)\le-\delta_{\cK}$ give
\eqref{eq:linearATsign} by uniform continuity.  Compactness gives
uniform strictness away from $q$.  This completes the proof.
\end{proof}

\subsection{Two-replica Guerra--Talagrand coercivity}\label{subsec:GT}

For an attainable overlap $v\in\mathcal R_N
\coloneqq\{-1,-1+2/N,\ldots,1\}$, define
\begin{equation}
 Z^{(2)}_{N,s}(v)
 \coloneqq
 \sum_{\substack{\sigma^1,\sigma^2\in\Sigma_N\\R_{12}=v}}
 e^{H_{N,s}(\sigma^1)+H_{N,s}(\sigma^2)}
 \label{eq:constrainedZ}
\end{equation}
and retain the notation $P_s^*$ from \eqref{eq:RSpathvalue}.  For
$\lambda\in\mathbb R$, define
\begin{equation}
 f_\lambda(x_1,x_2)
 \coloneqq\log\left(\frac14\sum_{\varepsilon_1,\varepsilon_2=\pm1}
 \exp\{\varepsilon_1x_1+\varepsilon_2x_2
       +\lambda\varepsilon_1\varepsilon_2\}\right).
 \label{eq:GTterminal}
\end{equation}

For $(\beta,h,s)\in\cK\times[0,1]$, retain $x_q$ and $\xi_s$ from
\eqref{eq:xqxi}.  Given $v\in[-1,1]$, write
$r_v\coloneqq|v|$ and put $\iota_v\coloneqq1$ for $v\ge0$ and
$\iota_v\coloneqq-1$ for $v<0$.  Define
\begin{align}
 Q^v_u
 &\coloneqq
 \begin{cases}
  \begin{pmatrix}u&\iota_vu\\ \iota_vu&u\end{pmatrix},
       &0\le u\le r_v,\\[2mm]
  \begin{pmatrix}u&v\\v&u\end{pmatrix},
       &r_v\le u\le1,
 \end{cases}
 \label{eq:GTpath}\\
 \gamma_v(u)
 &\coloneqq
 \begin{cases}
  \frac12x_q(u),&0\le u<r_v,\\
  x_q(u),&r_v\le u\le1.
 \end{cases}
 \label{eq:halfparameter}
\end{align}
Both $Q^v$ and
\begin{equation}
 A^v_s(u)\coloneqq\frac{\dd}{\dd u}\xi_s'(Q^v_u)
 =
 \begin{cases}
 s\beta^2\begin{pmatrix}1&\iota_v\\\iota_v&1\end{pmatrix},&u<r_v,\\[2mm]
 s\beta^2I_2,&u>r_v
 \end{cases}
 \label{eq:GTspeed}
\end{equation}
are positive semidefinite.  Scalar functions are applied entrywise to
matrices.  Let $U^{\lambda,v}_s$ be the finite-semigroup solution of
\begin{equation}
 \partial_uU+\frac12\left(
   \langle A^v_s,D^2U\rangle
   +\gamma_v(u)\langle A^v_s\nabla U,\nabla U\rangle
 \right)=0
 \label{eq:2DParisiPDE}
\end{equation}
with terminal condition $f_\lambda$ from \eqref{eq:GTterminal}.  With
$Y_0=h+\beta\sqrt{(1-s)q}\,Z_0$, set
\begin{align}
 \mathcal L_s(v)
 &\coloneqq\frac12\int_0^1\gamma_v(u)
   \sum_{a,b=1}^2 A^v_{s,ab}(u)Q^v_{u,ab}\,\dd u,
 \label{eq:GTcorrectiondef}\\
 \mathfrak P_s(\lambda,v)
 &\coloneqq2\log2+\E U^{\lambda,v}_s(0,Y_0,Y_0)
   -\lambda v-\mathcal L_s(v).
 \label{eq:GTfunctional}
\end{align}
Below $r_v$, all four summands in \eqref{eq:GTcorrectiondef} equal
$s\beta^2u$; above $r_v$, only the two diagonal summands remain.  Hence
\begin{equation}
 \mathcal L_s(v)
 =s\beta^2\int_0^1x_q(u)u\,\dd u
 =\frac{s\beta^2}{2}(1-q^2),
 \label{eq:GTcorrection}
\end{equation}
independently of $v$.

The required finite-volume Guerra--Talagrand upper bound is the
specialization proved in Appendix~\ref{app:specialGT}; it is an instance
of the general two-dimensional bound of Talagrand and Chen
\cite{Talagrand,ChenGT}.  In the present notation, it gives
\begin{equation}
 \frac1N\E\log Z^{(2)}_{N,s}(v)
 \le \mathfrak P_s(\lambda,v),
 \qquad \lambda\in\mathbb R.
 \label{eq:GTfunctionalbound}
\end{equation}

\begin{lemma}
\label{lem:GTcontinuity}
On every compact set on which $\lambda$ is bounded, the maps
\[
 (\beta,h,s,\lambda,v)\longmapsto
 \mathfrak P_s(\lambda,v)-2P_s^*,
 \qquad
 (\beta,h,s,\lambda,v)\longmapsto
 \partial_\lambda\mathfrak P_s(\lambda,v)
\]
are jointly continuous for
$(\beta,h,s,v)\in\cK\times[0,1]\times[-1,1]$.
\end{lemma}

\begin{proof}
Split $[0,1]$ at $q$ and $|v|$.  On each resulting interval the
solution is obtained by one of the three Gaussian recursion operators
\[
 \mathcal T_mF(x)\coloneqq m^{-1}\log\E e^{mF(x+G)},
 \qquad m\in\{1/2,1\},
 \qquad
 \mathcal T_0F(x)\coloneqq\E F(x+G),
\]
where the covariance of $G$ is the integral of $A_s^v$ over that
interval.  These covariances depend continuously on the parameters.
At $v=0$, at $v=\pm q$, or when another interval collapses, the
corresponding covariance tends to zero and its operator tends to the
identity.  The same holds at $s=0$, since every $s\beta^2$ covariance
vanishes.  The terminal function has uniformly linear growth on
bounded $\lambda$-sets, so Gaussian exponential integrability and
dominated convergence prove continuity of the finite composition.
For the multiplier derivative, differentiate the terminal function
and propagate it through the same tilted Gaussian expectations; its
absolute value is at most one.  Dominated convergence applies again.
Finally, $q$, $Y_0$, \eqref{eq:GTcorrection}, and $P_s^*$ depend
continuously on the parameters.  This proves both assertions.
\end{proof}

\begin{lemma}[Uniform Taylor coercivity]
\label{lem:Taylorcoercivity}
Let $p$ range over a compact set, let $q_p\in[-1,1]$, and let
$H_p(\lambda,v)$ be continuous for $|\lambda|\le\lambda_0$ and
$-1\le v\le1$.  Suppose it is twice differentiable in $\lambda$ and
\[
 |\partial_{\lambda\lambda}H_p(\lambda,v)|\le M.
\]
Assume that, whenever $|v-q_p|\le\varepsilon$,
\[
 H_p(0,v)\le0,\qquad
 \partial_\lambda H_p(0,v)(v-q_p)
 \le-c|v-q_p|^2,
 \qquad
 |\partial_\lambda H_p(0,v)|\le M\lambda_0,
\]
and that
\[
 \inf_{|\lambda|\le\lambda_0}H_p(\lambda,v)\le-\kappa
 \qquad\text{if }|v-q_p|\ge\varepsilon.
\]
Then, for a constant $c'>0$ independent of $p$,
\[
 \inf_{|\lambda|\le\lambda_0}H_p(\lambda,v)
 \le-c'|v-q_p|^2.
\]
\end{lemma}

\begin{proof}
In the local region, set
$\lambda=-\partial_\lambda H_p(0,v)/M$.  Taylor's theorem gives
\[
 H_p(\lambda,v)
 \le-\frac{|\partial_\lambda H_p(0,v)|^2}{2M}
 \le-\frac{c^2}{2M}|v-q_p|^2.
\]
Outside that region use the fixed loss and $|v-q_p|\le2$.
\end{proof}

\begin{proposition}
\label{prop:GTcoercivity}
There is a constant $c_{\cK}>0$ such that
\begin{equation}
 \frac1N\E\log Z^{(2)}_{N,s}(v)
 \le 2P_s^*-c_{\cK}(v-q)^2
 \label{eq:GTcoercivity}
\end{equation}
for all $(\beta,h,s)\in\cK\times[0,1]$, all $N$, and all
$v\in\mathcal R_N$.
\end{proposition}

\begin{proof}
We give the normalization in full.  This is important here: with a
different factor in either the two-dimensional cumulative parameter
or the Lagrange term, the mixed derivative below would not be the AT
replicon.

Recall $\xi_s$ from \eqref{eq:xqxi}.
The covariance of the centered Gaussian part of \eqref{eq:path} is
\begin{equation}
 \E H_{N,s}^{\rm cen}(\sigma)H_{N,s}^{\rm cen}(\tau)
 =N\xi_s(R(\sigma,\tau))-\frac{s\beta^2}{2}.
 \label{eq:GTcovariance}
\end{equation}
If $G\sim N(0,s\beta^2/2)$ is independent and is added to every
configuration, the covariance becomes exactly $N\xi_s(R)$.  The
one-replica log partition function is shifted by $G$, and the
two-replica log partition function by $2G$; their disorder
expectations are unchanged.  Define
\[
 \widetilde H_{N,s}^{\rm cen}(\sigma)
 \coloneqq H_{N,s}^{\rm cen}(\sigma)+G.
\]
For the remainder of this application we denote
$\widetilde H_{N,s}^{\rm cen}$ again by $H_{N,s}^{\rm cen}$.  Thus its
covariance is exactly $N\xi_s(R)$, and the normalization used in
Appendix~\ref{app:specialGT} introduces no finite-$N$ error.

After the harmless covariance normalization above, the specialized
Guerra--Talagrand bound from Appendix~\ref{app:specialGT} gives
\eqref{eq:GTfunctionalbound}.
The statement is finite-volume and includes the rank-one pieces of the
path; no limiting argument in $N$ is involved.

We next identify the zero-source derivatives.  Suppose first that
$v\ge0$.  On $[v,1]$, the PDE at $\lambda=0$ factorizes as
\[
 U^{0,v}_s(u,x_1,x_2)=\Psi(u,x_1)+\Psi(u,x_2).
\]
On $[0,v]$, set $V(u,x)\coloneqq U^{0,v}_s(u,x,x)$.  From
\eqref{eq:2DParisiPDE},
\[
 \partial_uV+\frac{s\beta^2}{2}
 \left(V_{xx}+\frac{x_q(u)}2V_x^2\right)=0,
 \qquad V(v,x)=2\Psi(v,x).
\]
Hence $V=2\Psi$.  Combining this identity with
\eqref{eq:GTcorrection} and the scalar trial value
\eqref{eq:scalartrial} gives the exact flatness identity
\begin{equation}
 \mathfrak P_s(0,v)=2P_s^*,
 \qquad 0\le v\le1.
 \label{eq:GTflat}
\end{equation}
We now derive the multiplier derivative explicitly.  Set
\[
 W^{v}_s(u,x_1,x_2)
 \coloneqq\left.\partial_\lambda
 U^{\lambda,v}_s(u,x_1,x_2)\right|_{\lambda=0}.
\]
Direct differentiation of the terminal function gives
\begin{equation}
 f_0(x_1,x_2)=\log\cosh x_1+\log\cosh x_2,\qquad
 \left.\partial_\lambda f_\lambda(x_1,x_2)\right|_{\lambda=0}
 =\tanh x_1\tanh x_2.
 \label{eq:GTterminalderivative}
\end{equation}
On $[v,1]$, $A_s^v=s\beta^2I_2$ and
$U_s^{0,v}=\Psi(x_1)+\Psi(x_2)$.  Differentiating
\eqref{eq:2DParisiPDE} gives
\begin{align}
 \partial_uW_s^v+\frac{s\beta^2}{2}\bigl\{
 \partial_{11}W_s^v+\partial_{22}W_s^v
 +2x_q(u)[\Psi_x(u,x_1)\partial_1W_s^v
          +\Psi_x(u,x_2)\partial_2W_s^v]\bigr\}=0.
 \label{eq:GTlinearizedupper}
\end{align}
Differentiating the scalar PDE in $x$ verifies that the solution with
terminal data \eqref{eq:GTterminalderivative} is
\begin{equation}
 W_s^v(u,x_1,x_2)
 =\Psi_x(u,x_1)\Psi_x(u,x_2),\qquad v\le u\le1.
 \label{eq:GTlinearizedproduct}
\end{equation}
On $[0,v]$, put $\overline W(u,x)\coloneqq W_s^v(u,x,x)$.  Since
$A_s^v=s\beta^2\begin{psmallmatrix}1&1\\1&1\end{psmallmatrix}$ and
$U_s^{0,v}(u,x,x)=2\Psi(u,x)$, the restricted linearized equation is
\begin{equation}
 \partial_u\overline W+\frac{s\beta^2}{2}\overline W_{xx}
 +s\beta^2x_q(u)\Psi_x(u,x)\overline W_x=0,\qquad
 \overline W(v,x)=\Psi_x(v,x)^2.
 \label{eq:GTlinearizedlower}
\end{equation}
The backward Kolmogorov formula for \eqref{eq:localfield} yields
\[
 \E W_s^v(0,Y_0,Y_0)=\E\Psi_x(v,X_v)^2=g_s(v).
\]
The correction $\mathcal L_s(v)$ is independent of $\lambda$, while
$\partial_\lambda(-\lambda v)=-v$.  We have therefore proved
\begin{equation}
 \partial_\lambda\mathfrak P_s(0,v)=g_s(v)-v.
 \label{eq:GTfirstderivative}
\end{equation}
In particular,
\begin{equation}
 \partial_v\partial_\lambda\mathfrak P_s(0,q)
 =s\beta^2\E[\partial_{xx}\Psi(q,X_q)^2]-1
 =s\alpha-1.
 \label{eq:GTmixedderivative}
\end{equation}
This is the replicon mixed derivative with all normalizing factors
displayed.

The differentiated PDE also gives a constant $M_{\cK}>0$ such that
\begin{equation}
 \sup_{\substack{(\beta,h,s)\in\cK\times[0,1]\\
                  -1\le v\le1,\ \lambda\in\mathbb R}}
 |\partial_{\lambda\lambda}\mathfrak P_s(\lambda,v)|
 \le M_{\cK}.
 \label{eq:GTsecondderivative}
\end{equation}
This bound follows directly from the finite-step semigroup formula.
For a Gaussian vector $G$ and $m>0$, write
\[
 \mathcal T_mF(x)=m^{-1}\log\E\exp\{mF(x+G)\},
 \qquad \mathcal T_0F(x)=\E F(x+G).
\]
Then
\begin{align*}
 \partial_\lambda\mathcal T_mF
 &=\E_m[\partial_\lambda F],\\
 \partial_{\lambda\lambda}\mathcal T_mF
 &=\E_m[\partial_{\lambda\lambda}F]
   +m\Var_m(\partial_\lambda F),
\end{align*}
where $\E_m$ denotes the exponentially tilted expectation.  At the
terminal time, $|\partial_\lambda f_\lambda|\le1$ and
$0\le\partial_{\lambda\lambda}f_\lambda\le1$.  For every signed path
$-1\le v\le1$, the only positive semigroup parameters are $1/2$ and
$1$; splitting at $q$ and $|v|$ therefore creates at most two
nonzero levels.  The zero level is an ordinary Gaussian expectation.
Iterating the displayed identities gives
$|\partial_\lambda U|\le1$ and, for example,
$0\le\partial_{\lambda\lambda}U\le3$.  Averaging the base Gaussian and
differentiating $-\lambda v$ do not change the second derivative.
This proves \eqref{eq:GTsecondderivative} on the full signed range,
uniformly in all parameters.

We next compute the signed derivative by the same semigroups.  Suppose
$-q\le v\le q$, put $r\coloneqq|v|$ and
$\iota\coloneqq\operatorname{sgn}(v)$, with $\iota\coloneqq1$ when
$v=0$.  The
rank-one interval $[0,r]$ lies below the first mass $q$.  At
$\lambda=0$, the $m=1$ recursion on $[q,1]$ and the diagonal heat
recursion on $[r,q]$ give
\[
 U_s^{0,v}(r,x_1,x_2)=\Psi(r,x_1)+\Psi(r,x_2).
\]
The remaining $m=0$ recursion is linear.  Starting it from
$(Y_0,Y_0)$ gives the pair
\[
 X_1\coloneqq Y_0+\beta\sqrt{s}\sqrt r\,Z,\qquad
 X_2\coloneqq Y_0+\iota\beta\sqrt{s}\sqrt r\,Z
\]
at time $r$.  Each marginal is the scalar heat evolution, so
\[
 \E U_s^{0,v}(0,Y_0,Y_0)=2\E\Psi(0,Y_0).
\]
Together with \eqref{eq:GTcorrection}, this proves the flatness
identity in \eqref{eq:signedflat}.

For the derivative, \eqref{eq:GTterminalderivative} and
\eqref{eq:GTlinearizedproduct} give
$W_s^v(r,x_1,x_2)=\Psi_x(r,x_1)\Psi_x(r,x_2)$.  The explicit scalar
heat semigroup gives
\[
 \Psi_x(r,x)=\E_G\tanh(x+\beta\sqrt{s}\sqrt{q-r}\,G).
\]
Introduce independent $G_1,G_2$, also independent of $(Y_0,Z)$, and
put
\[
 Y_1\coloneqq Y_0+\beta\sqrt{s}\sqrt r\,Z
 +\beta\sqrt{s}\sqrt{q-r}\,G_1,\qquad
 Y_2\coloneqq Y_0+\iota\beta\sqrt{s}\sqrt r\,Z
 +\beta\sqrt{s}\sqrt{q-r}\,G_2.
\]
Then both means are $h$, both variances are
$\beta^2(1-s)q+s\beta^2q=\beta^2q$, and their covariance is
$\beta^2(1-s)q+\iota s\beta^2r
=\beta^2(1-s)q+s\beta^2v$.  It follows that
\[
 \E W_s^v(0,Y_0,Y_0)=\E\tanh(Y_1)\tanh(Y_2)=f_s(v).
\]
After differentiating $-\lambda v$, we have proved, rather than
assumed,
\begin{equation}
 \mathfrak P_s(0,v)=2P_s^*,\qquad
 \partial_\lambda\mathfrak P_s(0,v)=f_s(v)-v,
 \label{eq:signedflat}
\end{equation}
where
\begin{equation}
 f_s(v)\coloneqq\E\tanh(Y_1(v))\tanh(Y_2(v)),
 \quad
 \begin{cases}
 \operatorname{Var}Y_1(v)=\operatorname{Var}Y_2(v)=\beta^2q,\\
 \operatorname{Cov}(Y_1(v),Y_2(v))
 =\beta^2(1-s)q+s\beta^2v,
 \end{cases}
 \label{eq:signedf}
\end{equation}
and both Gaussian variables have mean $h$.  Gaussian differentiation
and Cauchy--Schwarz give
\begin{equation}
 f_s'(v)=s\beta^2\E[\sech^2Y_1(v)\sech^2Y_2(v)]
 \le s\alpha.
 \label{eq:signedcontraction}
\end{equation}
Since $f_s(q)=q$, it follows that
\begin{equation}
 f_s(v)-v\ge(1-s\alpha)(q-v)
 \qquad(-q\le v<q).
 \label{eq:signedslope}
\end{equation}
Together with \eqref{eq:GTfirstderivative}, define on $[-q,1]$
\begin{equation}
 \eta_s(v)\coloneqq\partial_\lambda\mathfrak P_s(0,v)
 =
 \begin{cases}
  f_s(v)-v,&-q\le v\le q,\\
  g_s(v)-v,&q\le v\le1.
 \end{cases}
 \label{eq:GTeta}
\end{equation}
Equations \eqref{eq:signedslope} and \eqref{eq:linearATsign} show that,
in a uniform neighborhood of $q$,
\begin{equation}
 \eta_s(v)(v-q)\le-c_{\cK}|v-q|^2.
 \label{eq:GTetaslope}
\end{equation}

For $v<-q$, the strict signed gap can be computed directly.  Write
$r\coloneqq|v|$ and $t\coloneqq s\beta^2(r-q)$.  When $s>0$, one has
$t>0$.
On $[r,1]$
the zero-source PDE factorizes, and since
\[
 \Psi(u,x)=\log\cosh x+\frac{s\beta^2}{2}(1-u),
 \qquad q\le u\le1,
\]
its value at time $r$ is
\[
 U_s^{0,v}(r,x_1,x_2)
 =\log\cosh x_1+\log\cosh x_2+s\beta^2(1-r).
\]
On $[q,r]$,
$A_s^v=s\beta^2\begin{psmallmatrix}1&-1\\-1&1\end{psmallmatrix}$
and $\gamma_v=1/2$.  Thus $W\coloneqq\exp(U_s^{0,v}/2)$ solves the linear
backward heat equation in the direction $(1,-1)$, and
\begin{align}
 U_s^{0,v}(q,x_1,x_2)
 ={}&s\beta^2(1-r)\nonumber\\
 &+2\log\E\sqrt{
   \cosh(x_1+\sqrt t\,Z)\cosh(x_2-\sqrt t\,Z)}.
 \label{eq:signedhopfcole}
\end{align}
Cauchy--Schwarz and
$\E\cosh(x+\sqrt t\,Z)=e^{t/2}\cosh x$ give
\begin{equation}
 U_s^{0,v}(q,x_1,x_2)
 \le \Psi(q,x_1)+\Psi(q,x_2).
 \label{eq:signedCS}
\end{equation}
Equality in Cauchy--Schwarz would require
$\cosh(x_1+y)$ and $\cosh(x_2-y)$ to be proportional for every
$y\in\mathbb R$, which is equivalent to $x_1=-x_2$.

Below $q$, the PDE is linear.  Starting from $(Y_0,Y_0)$, its two
coordinates at time $q$ are
\[
 (Y_0+\beta\sqrt{s}\sqrt q\,Z,\;
  Y_0-\beta\sqrt{s}\sqrt q\,Z).
\]
Their sum is $2Y_0$.  Since $h>0$, one has
$\mathbb P(Y_0=0)=0$ (also when $s=1$, because then $Y_0=h$).
Therefore \eqref{eq:signedCS} is strict after averaging.  Its two
marginals both have the scalar local-field law, so
\begin{equation}
 \E U_s^{0,v}(0,Y_0,Y_0)<2\E\Psi(0,Y_0),
 \qquad v<-q, s>0.
 \label{eq:signedstrictPDE}
\end{equation}
The correction identity \eqref{eq:GTcorrection} now yields
\begin{equation}
 \mathfrak P_s(0,v)<2P_s^*,
 \qquad v<-q, s>0.
 \label{eq:signedpositivity}
\end{equation}
This is the signed GT positivity step in the present RS
normalization, derived without a Parisi free-energy formula.

We now prove the required uniformity, paying particular attention to
$s\downarrow0$.  At $s=0$ the interaction coefficient
$\beta\sqrt{s}$ vanishes, and
\[
 \mathfrak P_0(\lambda,v)
 =2\log2+\E f_\lambda(Y_0,Y_0)-\lambda v.
\]
The first derivative at zero is $q-v$, and the second derivative is
the variance, under the one-site coupled measure, of
$\varepsilon_1\varepsilon_2$; it is at most one.  Taylor's theorem
therefore gives, for every $\lambda\in\mathbb R$,
\begin{equation}
 \mathfrak P_0(\lambda,v)
 \le2P_0^*+\lambda(q-v)+\frac{\lambda^2}{2}.
 \label{eq:GTendpoint}
\end{equation}
For $v\le-q$, take $\lambda=v-q$.  By \eqref{eq:qcompact},
\begin{equation}
 \mathfrak P_0(v-q,v)
 \le2P_0^*-\frac12(v-q)^2
 \le2P_0^*-2q_{\cK}^2.
 \label{eq:GTsmallsgapzero}
\end{equation}
By Lemma~\ref{lem:GTcontinuity},
$\mathfrak P_s(\lambda,v)-2P_s^*$ is jointly continuous in
$(\beta,h,s,\lambda,v)$ on compact sets, including at $s=0$, $v=0$,
and $v=\pm q$.  Applying this continuity to
$\lambda=v-q$ and the compact set
\[
 \{(\beta,h,v):(\beta,h)\in\cK,\ -1\le v\le-q(\beta,h)\}
\]
shows that there exists $s_0>0$ such that
\begin{equation}
 \mathfrak P_s(v-q,v)\le2P_s^*-q_{\cK}^2,
 \qquad 0\le s\le s_0,\quad v\le-q.
 \label{eq:GTsmallsgap}
\end{equation}

It remains to treat $s\in[s_0,1]$.  At $v=-q$,
\eqref{eq:signedslope} gives
\[
 \partial_\lambda\mathfrak P_s(0,-q)
 \ge2\delta_{\cK}q_{\cK}.
\]
Joint continuity of this derivative from Lemma~\ref{lem:GTcontinuity} and
\eqref{eq:GTsecondderivative} yield $\eta_{\cK}>0$ such that, for
$-q-\eta_{\cK}\le v\le-q$ and $s\in[s_0,1]$,
\[
 \partial_\lambda\mathfrak P_s(0,v)
 \ge\delta_{\cK}q_{\cK}.
\]
The Cauchy--Schwarz argument above gives
$\mathfrak P_s(0,v)\le2P_s^*$ throughout this strip.  Taylor's theorem,
with
\[
 \ell_{\cK}\coloneqq
 \frac{\delta_{\cK}q_{\cK}}{M_{\cK}},
 \qquad\lambda\coloneqq-\ell_{\cK},
\]
therefore gives
\[
 \mathfrak P_s(-\ell_{\cK},v)
 \le2P_s^*-\frac12\delta_{\cK}q_{\cK}\ell_{\cK}.
\]
On the remaining compact set
\[
 \{(\beta,h,s,v):(\beta,h)\in\cK,\ s_0\le s\le1,\
                 -1\le v\le-q-\eta_{\cK}\},
\]
\eqref{eq:signedpositivity} is strict and continuous, so its minimum
gap is also positive.  Combining this with
\eqref{eq:GTsmallsgap} proves that, for some
$\kappa_{\cK}>0$,
\begin{equation}
 \inf_{\lambda\in\mathbb R}\mathfrak P_s(\lambda,v)
 \le2P_s^*-\kappa_{\cK},\qquad v\le-q,
 \label{eq:GTuniformnegativegap}
\end{equation}
uniformly on $\cK\times[0,1]$.

We finish with one Taylor argument.  Put
\[
 H_{(\beta,h,s)}(\lambda,v)
 \coloneqq\mathfrak P_s(\lambda,v)-2P_s^*.
\]
In a uniform neighborhood of $q$, the flatness identities
\eqref{eq:GTflat} and \eqref{eq:signedflat}, together with
\eqref{eq:GTetaslope}, verify the local hypotheses of
Lemma~\ref{lem:Taylorcoercivity}; shrink the neighborhood if necessary.
Outside that neighborhood, \eqref{eq:signedslope} on $[-q,q)$ and
\eqref{eq:ATsign} on $(q,1]$ give a fixed loss by a bounded multiplier,
using \eqref{eq:GTsecondderivative} and compactness.  The range
$v\le-q$ is covered by \eqref{eq:GTuniformnegativegap}; all the
multipliers used in its proof lie in a fixed compact interval.
Lemma~\ref{lem:Taylorcoercivity} therefore yields
\[
 \inf_{\lambda\in\mathbb R}\mathfrak P_s(\lambda,v)
 \le2P_s^*-c_{\cK}(v-q)^2.
\]
Combining this with \eqref{eq:GTfunctionalbound} proves
\eqref{eq:GTcoercivity}.
\end{proof}

\begin{remark}
The additional parameter \eqref{eq:halfparameter} is the missing
parameter in the scalar Lata\l a interpolation.  If it is omitted, the
mixed derivative contains anomalous and longitudinal fluctuations.
With it, \eqref{eq:GTmixedderivative} is exactly
$-(1-s\alpha)$.
\end{remark}

\section{Preliminary overlap concentration}\label{sec:concentration}

\subsection{Finite-volume Gronwall closure}\label{subsec:gronwall}

Choose once and for all
\begin{equation}
 0<\lambda_0<c_{\cK}.
 \label{eq:lambda0}
\end{equation}
Define the quadratically coupled pressure and its normalized excess by
\begin{align}
 p^{(2)}_{N,s}(\lambda_0)
 &\coloneqq
 \frac1{2N}\E\log
 \sum_{\sigma^1,\sigma^2}
 \exp\left\{
 H_{N,s}(\sigma^1)+H_{N,s}(\sigma^2)
 +\frac{\lambda_0N}{2}Q_{12}^2\right\},
 \label{eq:quadraticpressure}\\
 F_{N,s}(\lambda_0)
 &\coloneqq
 p^{(2)}_{N,s}(\lambda_0)-\phi_N(s)
 =
 \frac1{2N}\E\log
 \braket{e^{\lambda_0NQ_{12}^2/2}}_s.
 \label{eq:normalizedcoupling}
\end{align}

\begin{lemma}
\label{lem:sublinearpressure}
With
\[
 \epsilon_N\coloneqq C_{\cK}
 \sqrt{\frac{\log(N+1)}N}+\frac{C_{\cK}\log(N+1)}N,
\]
one has
\begin{equation}
 p^{(2)}_{N,s}(\lambda_0)\le P_s^*+\epsilon_N
 \label{eq:pressureenvelope}
\end{equation}
uniformly on $\cK\times[0,1]$.
\end{lemma}

\begin{proof}
Let $L_v\coloneqq\log Z^{(2)}_{N,s}(v)$.  Each $L_v$ is a Lipschitz function
of the Gaussian disorder with squared Lipschitz constant at most
$C_{\cK}N$.  Indeed,
\[
 \left|\partial_{g_{ij}}L_v\right|
 \le\frac{2\beta\sqrt s}{\sqrt N},\qquad
 \left|\partial_{z_i}L_v\right|
 \le2\beta\sqrt{(1-s)q},
\]
and summing the squares over $i<j$ and $i$ gives the claim.  Hence
\begin{equation}
 \E\max_{v\in\mathcal R_N}(L_v-\E L_v)
 \le C_{\cK}\sqrt{N\log(N+1)}.
 \label{eq:gaussianmax}
\end{equation}
Using $|\mathcal R_N|=N+1$, \eqref{eq:GTcoercivity}, and
$\lambda_0<c_{\cK}$,
\begin{align*}
 2Np^{(2)}_{N,s}(\lambda_0)
 &=
 \E\log\sum_{v\in\mathcal R_N}
 \exp\left\{L_v+\frac{\lambda_0N}{2}(v-q)^2\right\}\\
 &\le
 \log(N+1)+
 \max_v\left\{\E L_v+\frac{\lambda_0N}{2}(v-q)^2\right\}
 +C_{\cK}\sqrt{N\log(N+1)}\\
 &\le2NP_s^*+2N\epsilon_N.
\end{align*}
This proves \eqref{eq:pressureenvelope}.
\end{proof}

Put
\begin{equation}
 D_N(s)\coloneqq P_s^*-\phi_N(s).
 \label{eq:DN}
\end{equation}
At $s=0$, $D_N(0)=0$.  For $0<s<1$, differentiating under the
expectation and applying \eqref{eq:GIBP} to the variables $g_{ij}$ and
$z_i$ gives
\begin{align*}
 \phi_N'(s)
 ={}&\frac{\beta^2}{2N^2}\sum_{i<j}
 \E\left[1-\braket{\sigma_i\sigma_j}_s^2\right]
 -\frac{\beta^2q}{2N}\sum_i
 \E\left[1-\braket{\sigma_i}_s^2\right]\\
 ={}&\frac{\beta^2}{4}
 \left((1-q)^2-\nu_s[Q_{12}^2]\right).
\end{align*}
Here we used
\[
 \sum_{i<j}\sigma_i^1\sigma_j^1\sigma_i^2\sigma_j^2
 =\frac12(N^2R_{12}^2-N),
 \qquad
 \frac1N\sum_i\E\braket{\sigma_i^1\sigma_i^2}_s
 =\nu_s[R_{12}].
\]
Combining this identity with \eqref{eq:RSpathvalue} yields
\begin{equation}
 D_N'(s)=\frac{\beta^2}{4}A_s.
 \label{eq:DNprime}
\end{equation}
Both sides extend continuously to $s=0,1$ by finite-dimensional
Gaussian dominated convergence, so the identity also holds for the
one-sided endpoint derivatives.
On the other hand, convexity in the quadratic-coupling parameter gives
\[
 \frac14A_s
 =\left.\partial_\lambda F_{N,s}(\lambda)\right|_{\lambda=0+}
 \le\frac{F_{N,s}(\lambda_0)}{\lambda_0}.
\]
By \eqref{eq:pressureenvelope},
\[
 F_{N,s}(\lambda_0)
 \le D_N(s)+\epsilon_N.
\]
Consequently,
\begin{equation}
 D_N'(s)
 \le\frac{\beta^2}{\lambda_0}
 \bigl(D_N(s)+\epsilon_N\bigr).
 \label{eq:gronwallDN}
\end{equation}
Gronwall's inequality and compactness of $\cK$ imply
\begin{equation}
 \sup_{\cK\times[0,1]}D_N(s)\le C_{\cK}\epsilon_N,
 \qquad
 \sup_{\cK\times[0,1]}A_s\le C_{\cK}\epsilon_N.
 \label{eq:preconcentration}
\end{equation}
This is the first, deliberately nonoptimal, concentration estimate.
Its role is to normalize the coupled pressure without invoking the
Parisi formula.

\subsection{Fixed-deviation overlap concentration}\label{subsec:tails}

\begin{corollary}
\label{cor:tail}
For every $\varepsilon>0$, there are constants
$c_{\cK,\varepsilon},C_{\cK,\varepsilon}>0$ such that
\begin{equation}
 \sup_{\substack{(\beta,h)\in\cK\\0\le s\le1}}
 \nu_s\bigl[|Q_{12}|\ge\varepsilon\bigr]
 \le C_{\cK,\varepsilon}e^{-c_{\cK,\varepsilon}N}.
 \label{eq:tail}
\end{equation}
\end{corollary}

\begin{proof}
Let
\[
 Y_{N,s}
 \coloneqq\log\braket{e^{\lambda_0NQ_{12}^2/2}}_s.
\]
By \eqref{eq:normalizedcoupling},
\eqref{eq:pressureenvelope}, and
\eqref{eq:preconcentration},
\[
 \E Y_{N,s}=2NF_{N,s}(\lambda_0)
 \le C_{\cK}N\epsilon_N=o(N)
\]
uniformly.  As a difference of one coupled and two uncoupled log
partition functions, $Y_{N,s}$ has squared Gaussian Lipschitz constant
at most $C_{\cK}N$.  Gaussian concentration therefore gives, for all
large enough $N$,
\[
 \mathbb P\left(
 Y_{N,s}>\frac{\lambda_0\varepsilon^2N}{4}\right)
 \le C_{\cK,\varepsilon}e^{-c_{\cK,\varepsilon}N}.
\]
Outside this exceptional event,
\[
 \braket{\1_{\{|Q_{12}|\ge\varepsilon\}}}_s
 \le
 \exp\left\{-\frac{\lambda_0\varepsilon^2N}{2}
 +Y_{N,s}\right\}
 \le e^{-\lambda_0\varepsilon^2N/4}.
\]
Taking disorder expectation and enlarging the constant for the
finitely many remaining $N$ proves \eqref{eq:tail}.
\end{proof}

\section{Cavity closure}\label{sec:cavityclosure}

\subsection{The finite-size cavity system}\label{subsec:cavity}

The next proposition is the one-species cavity fluctuation system.  We
include its derivation because the size and form of its remainder are
essential for the absorption step.
Recall from \eqref{eq:ABC} that
\[
 A_s=\nu_s[Q_{12}^2],\qquad
 B_s=\nu_s[Q_{12}Q_{13}],\qquad
 C_s=\nu_s[Q_{12}Q_{34}].
\]

\begin{proposition}
\label{prop:cavity}
Let $x_s\coloneqq(A_s,B_s,C_s)^{\mathsf T}$, and define
$\rho_{N,s}\in\R^3$ by
\begin{equation}
 (I-sL)x_s=\frac1N\theta+\rho_{N,s},
 \label{eq:cavitysystem}
\end{equation}
where
\begin{equation}
 \theta\coloneqq
 \begin{pmatrix}
 1-q^2\\ q-q^2\\ r-q^2
 \end{pmatrix},
 \label{eq:theta}
\end{equation}
\begin{equation}
 L\coloneqq
 \begin{pmatrix}
 b_2&-4b_1&3b_0\\
 b_1&b_2-2b_1-3b_0&6b_0-3b_1\\
 b_0&4b_1-8b_0&b_2-8b_1+10b_0
 \end{pmatrix},
 \label{eq:L}
\end{equation}
with
\begin{equation}
 b_2\coloneqq\beta^2(1-q^2),\qquad
 b_1\coloneqq\beta^2q(1-q),\qquad
 b_0\coloneqq\beta^2(r-q^2),
 \label{eq:bs}
\end{equation}
Then, uniformly on $\cK\times[0,1]$,
\begin{equation}
 \norm{\rho_{N,s}}_\infty
 \le C_{\cK}
 \left(N^{-3/2}+\nu_s[|Q_{12}|^3]\right).
 \label{eq:rho}
\end{equation}
\end{proposition}

\begin{proof}
Write $\varepsilon_a\coloneqq\sigma_N^a$ and
\[
 Q^-_{ab}\coloneqq\frac1N\sum_{i<N}\sigma_i^a\sigma_i^b-q.
\]
By site exchangeability after disorder averaging,
\begin{align}
 A_s
 &=\nu_s\bigl[Q_{12}(\varepsilon_1\varepsilon_2-q)\bigr],\nonumber\\
 B_s
 &=\nu_s\bigl[Q_{12}(\varepsilon_1\varepsilon_3-q)\bigr],\nonumber\\
 C_s
 &=\nu_s\bigl[Q_{12}(\varepsilon_3\varepsilon_4-q)\bigr].
 \label{eq:siteexchange}
\end{align}
Since
$Q_{12}=Q^-_{12}+N^{-1}\varepsilon_1\varepsilon_2$,
each line consists of a cavity term involving $Q^-_{12}$ and a
diagonal term of order $N^{-1}$; explicitly,
\begin{align}
 A_s&=\nu_s[Q^-_{12}(\varepsilon_1\varepsilon_2-q)]
 +\frac1N\nu_s[\varepsilon_1\varepsilon_2
                    (\varepsilon_1\varepsilon_2-q)],\nonumber\\
 B_s&=\nu_s[Q^-_{12}(\varepsilon_1\varepsilon_3-q)]
 +\frac1N\nu_s[\varepsilon_1\varepsilon_2
                    (\varepsilon_1\varepsilon_3-q)],\nonumber\\
 C_s&=\nu_s[Q^-_{12}(\varepsilon_3\varepsilon_4-q)]
 +\frac1N\nu_s[\varepsilon_1\varepsilon_2
                    (\varepsilon_3\varepsilon_4-q)].
 \label{eq:cavitydecomposition}
\end{align}

We define the last-spin interpolation explicitly.  Write
$\sigma\coloneqq(\rho,\varepsilon)$ with
$\rho\in\{-1,1\}^{N-1}$ and let
\begin{align*}
 H^-_{N,s}(\rho)
 \coloneqq{}&\frac{\beta\sqrt s}{\sqrt N}
 \sum_{1\le i<j<N}g_{ij}\rho_i\rho_j
 +\sum_{i<N}\bigl(h+\beta\sqrt{(1-s)q}\,z_i\bigr)\rho_i.
\end{align*}
Let $\widehat z$ be a standard Gaussian independent of all existing
disorder.  For $0\le u\le1$, set
\begin{align}
 H_{N,s,u}(\rho,\varepsilon)
 \coloneqq{}&H^-_{N,s}(\rho)
 +\bigl(h+\beta\sqrt{(1-s)q}\,z_N\bigr)\varepsilon\nonumber\\
 &+\frac{\beta\sqrt{su}}{\sqrt N}
   \sum_{i<N}g_{iN}\rho_i\varepsilon
 +\beta\sqrt{s(1-u)q}\,\widehat z\,\varepsilon .
 \label{eq:cavityHamiltonian}
\end{align}
Denote its disorder--Gibbs expectation by $\nu_{s,u}$.  At $u=1$,
\eqref{eq:cavityHamiltonian} is exactly \eqref{eq:path}, so
$\nu_{s,1}=\nu_s$.  At $u=0$, the last spin is independent of the
first $N-1$ spins and sees the Gaussian field
\[
 Y\stackrel d=h+\beta\sqrt q\,Z.
\]
Conditional on $Y$, its replicas are independent, with mean
$\tanh Y$.  Thus, at the decoupled endpoint,
\begin{equation}
 \E_0[\varepsilon_a\varepsilon_b]=q,\qquad
 \E_0[\varepsilon_a\varepsilon_b\varepsilon_c\varepsilon_d]=r
 \label{eq:epsilonmoments}
\end{equation}
for distinct indices.

For completeness, we differentiate this interpolation.  The
derivative of the covariance between the last-spin parts of two
configurations $(\rho,\varepsilon)$ and
$(\rho',\varepsilon')$ is
\begin{equation}
 s\beta^2\varepsilon\varepsilon'
 \left(\frac1N\sum_{i<N}\rho_i\rho_i'-q\right)
 =s\beta^2\varepsilon\varepsilon'Q^-(\rho,\rho').
 \label{eq:cavitycovariancederivative}
\end{equation}
For a bounded function $F$ of $n$ replicas, differentiating the
normalized Gibbs expectation for $0<u<1$ and applying Gaussian
integration by parts to \eqref{eq:cavitycovariancederivative} gives
\begin{align}
 \frac{\dd}{\dd u}\nu_{s,u}[F]
 =s\beta^2\Bigg\{&
 \sum_{1\le a<b\le n}
 \nu_{s,u}[F\varepsilon_a\varepsilon_bQ^-_{ab}]
 -n\sum_{a=1}^n
 \nu_{s,u}[F\varepsilon_a\varepsilon_{n+1}Q^-_{a,n+1}]
 \nonumber\\
 &+\frac{n(n+1)}2
 \nu_{s,u}[F\varepsilon_{n+1}\varepsilon_{n+2}
 Q^-_{n+1,n+2}]\Bigg\}.
 \label{eq:cavityderivative}
\end{align}
Both sides of \eqref{eq:cavityderivative} extend continuously to
$u=0,1$.  We use those extensions for the one-sided endpoint
derivatives below; thus the square-root coefficients in
\eqref{eq:cavityHamiltonian} cause no endpoint singularity.
To see all coefficients in one formula, define the replica operator
\begin{align*}
 \mathcal D_nF\coloneqq{}&
 \sum_{1\le a<b\le n}
 F\varepsilon_a\varepsilon_bQ^-_{ab}
 -n\sum_{a=1}^n
 F\varepsilon_a\varepsilon_{n+1}Q^-_{a,n+1}\\
 &+\frac{n(n+1)}2
 F\varepsilon_{n+1}\varepsilon_{n+2}Q^-_{n+1,n+2}.
\end{align*}
Then \eqref{eq:cavityderivative} is simply
$\frac{\dd}{\dd u}\nu_{s,u}[F]
 =s\beta^2\nu_{s,u}[\mathcal D_nF]$.  This form also makes the
second differentiation below unambiguous: each application of
$\mathcal D$ appends exactly one centered cavity-overlap factor, and
the replica number is increased only to include the displayed new
indices.

\medskip
\noindent\emph{Cavity replacement lemma.}
The following estimate permits us to replace every intermediate
quadratic cavity moment by its full endpoint counterpart.  By
construction, $\nu_{s,1}=\nu_s$.  Set
\begin{equation}
 M_3(u)\coloneqq\nu_{s,u}[|Q^-_{12}|^3].
 \label{eq:cavityM3}
\end{equation}
Applying \eqref{eq:cavityderivative} with
$F=|Q^-_{12}|^3$, using $|Q^-_{ab}|\le2$ and replica symmetry, gives
\begin{equation}
 |M_3'(u)|\le C_{\cK}M_3(u).
 \label{eq:cavityM3derivative}
\end{equation}
Indeed, every term on the right is bounded by
$2\nu_{s,u}[|Q^-_{12}|^3]$.  Integrating backward from $u=1$ and applying
Gronwall's inequality yields
\begin{equation}
 \sup_{0\le u\le1}M_3(u)\le C_{\cK}M_3(1).
 \label{eq:cavityM3comparison}
\end{equation}
Since
$Q^-_{12}=Q_{12}-N^{-1}\varepsilon_1\varepsilon_2$ and
$|\varepsilon_1\varepsilon_2|=1$, the inequality
$(x+y)^3\le4(x^3+y^3)$ for $x,y\ge0$ gives
\begin{equation}
 M_3(1)\le C\left(\nu_s[|Q_{12}|^3]+N^{-3}\right).
 \label{eq:cavityM3endpoint}
\end{equation}

Next let $P^-$ be any product of two centered cavity overlaps, for
example $(Q^-_{12})^2$, $Q^-_{12}Q^-_{13}$, or
$Q^-_{12}Q^-_{34}$.  Differentiating $\nu_{s,u}[P^-]$ by
\eqref{eq:cavityderivative} produces products of three centered cavity
overlaps, possibly multiplied by spin variables.  By H\"older's
inequality and replica symmetry, each such term is bounded by $M_3(u)$.
Consequently,
\begin{equation}
 \left|\nu_{s,u}[P^-]-\nu_{s,v}[P^-]\right|
 \le C_{\cK}|u-v|M_3(1),
 \qquad 0\le u,v\le1.
 \label{eq:cavityquadraticcomparison}
\end{equation}
At the full endpoint, replacing either factor $Q^-_{ab}$ by $Q_{ab}$
costs at most
$C(N^{-1}|Q_{cd}|+N^{-2})$.  The pointwise Young inequality
\begin{equation}
 \frac{|x|}{N}\le\frac{|x|^3}{3}+\frac{2}{3N^{3/2}}
 \label{eq:cavityYoung}
\end{equation}
therefore implies, for each of the three quadratic overlap patterns,
\begin{equation}
 \left|\nu_{s,0}[P^-]-\nu_s[P]\right|
 \le C_{\cK}\left(\nu_s[|Q_{12}|^3]+N^{-3/2}\right),
 \label{eq:cavityendpointcomparison}
\end{equation}
where $P$ denotes the corresponding product of full centered overlaps.

\medskip
\noindent\emph{Cavity Taylor lemma.}
For each of the three centered cavity functions $F$ in
\eqref{eq:cavitydecomposition}, one has $\nu_{s,0}[F]=0$ and
\begin{equation}
 \nu_{s,1}[F]=
 \left.\frac{\dd}{\dd u}\nu_{s,u}[F]\right|_{u=0}
 +O_{\cK}\left(\nu_s[|Q_{12}|^3]+N^{-3/2}\right).
 \label{eq:cavityTaylorlemma}
\end{equation}
Indeed, the vanishing follows from \eqref{eq:epsilonmoments}, and the
first-order Taylor formula gives the derivative in
\eqref{eq:cavityTaylorlemma} plus
\[
 \int_0^1(1-u)\frac{\dd^2}{\dd u^2}\nu_{s,u}[F]\,\dd u.
\]
Formula \eqref{eq:cavityderivative} first
expresses the derivatives at $u=0$ in terms of the decoupled quadratic
moments.  Replacing those moments by $A_s,B_s,C_s$ using
\eqref{eq:cavityendpointcomparison} changes each equation by at most
$C_{\cK}(\nu_s[|Q_{12}|^3]+N^{-3/2})$.  Consequently,
\begin{align}
 A_s={}&s\bigl(b_2A_s-4b_1B_s+3b_0C_s\bigr)
       +\frac{1-q^2}{N}
       +O_{\cK}\left(\nu_s[|Q_{12}|^3]+N^{-3/2}\right),\nonumber\\
 B_s={}&s\bigl(b_1A_s+(b_2-2b_1-3b_0)B_s
                    +(6b_0-3b_1)C_s\bigr)
       +\frac{q-q^2}{N}
       +O_{\cK}\left(\nu_s[|Q_{12}|^3]+N^{-3/2}\right),\nonumber\\
 C_s={}&s\bigl(b_0A_s+(4b_1-8b_0)B_s
                    +(b_2-8b_1+10b_0)C_s\bigr)
       +\frac{r-q^2}{N}
       +O_{\cK}\left(\nu_s[|Q_{12}|^3]+N^{-3/2}\right).
 \label{eq:cavitytable}
\end{align}
For example, the first line uses $n=2$.  Before replacing decoupled
quadratic moments by endpoint moments, the three terms in
\eqref{eq:cavityderivative} are
\[
 (1-q^2)\nu_{s,0}[(Q^-_{12})^2],\qquad
 -4q(1-q)\nu_{s,0}[Q^-_{12}Q^-_{13}],\qquad
 3(r-q^2)\nu_{s,0}[Q^-_{12}Q^-_{34}].
\]
Equation \eqref{eq:cavityendpointcomparison} turns these into the first
identity in \eqref{eq:cavitytable}, with an error of the required size.
Here is a direct rule that verifies every entry.  Conditional on $Y$,
the $\varepsilon_a$ are independent with mean $\tanh Y$.  Hence a
spin monomial has expectation $1$, $q$, or $r$ according as the
number of replica indices occurring an odd number of times is
$0$, $2$, or $4$.  In particular, if $e$ is a replica edge, then
\[
 \E_0[(\varepsilon_p\varepsilon_q-q)\varepsilon_e]
 =
 \begin{cases}
  1-q^2,&e=\{p,q\},\\
  q-q^2,&e\text{ shares exactly one endpoint with }\{p,q\},\\
  r-q^2,&e\text{ is disjoint from }\{p,q\}.
 \end{cases}
\]
Counting the old, one-new, and two-new edges in $\mathcal D_2$,
$\mathcal D_3$, and $\mathcal D_4$ gives respectively
\[
 (b_2,-4b_1,3b_0),
\]
\[
 (b_1,\ b_2-2b_1-3b_0,\ 6b_0-3b_1),
\]
and
\[
 (b_0,\ 4b_1-8b_0,\ b_2-8b_1+10b_0),
\]
which verifies all three identities in \eqref{eq:cavitytable}.

The diagonal terms come from
$N^{-1}\varepsilon_1\varepsilon_2$ in the exact identity
$Q_{12}=Q^-_{12}+N^{-1}\varepsilon_1\varepsilon_2$.  For example,
\[
 \E_0[\varepsilon_1\varepsilon_2
       (\varepsilon_1\varepsilon_3-q)]
 =\E\tanh^2Y-q^2=q-q^2.
\]
The other two diagonal expectations are
$\E_0[\varepsilon_1\varepsilon_2
(\varepsilon_1\varepsilon_2-q)]=1-q^2$ and
$\E_0[\varepsilon_1\varepsilon_2
(\varepsilon_3\varepsilon_4-q)]=r-q^2$.

It remains to bound the integral Taylor remainder.  Each of the three
cavity functions is a bounded spin factor times one centered cavity
overlap.  The first application of $\mathcal D_n$ gives a finite sum
of bounded spin factors times two centered overlaps.  Applying
\eqref{eq:cavityderivative} once more, now with all replica indices
appearing in each summand included, gives a finite sum of terms
\[
 c\,\nu_{s,u}\bigl[
 \chi(\varepsilon)\,Q^-_{a_1b_1}Q^-_{a_2b_2}Q^-_{a_3b_3}\bigr],
\]
where $|c|\le C$ and $|\chi|\le C$; the number of terms is bounded
because the original replica numbers are only $2$, $3$, and $4$.
H\"older's inequality,
replica symmetry, \eqref{eq:cavityM3comparison}, and
\eqref{eq:cavityM3endpoint} give, uniformly in $u$,
\begin{equation}
 \left|\frac{\dd^2}{\dd u^2}\nu_{s,u}[F]\right|
 \le C_{\cK}\left(\nu_s[|Q_{12}|^3]+N^{-3}\right)
 \label{eq:cavitysecondderivativebound}
\end{equation}
for each of the three cavity functions $F$ occurring in
\eqref{eq:cavitydecomposition}.  Hence their integral Taylor remainders obey
the same bound.

For the $N^{-1}$ diagonal terms, one differentiation gives a bounded
spin observable times one centered cavity overlap.  Thus their change
between $u=0$ and $u=1$ is bounded by
\[
 \frac{C_{\cK}}N\sup_{0\le u\le1}\nu_{s,u}[|Q^-_{12}|].
\]
Using \eqref{eq:cavityYoung} pointwise and then
\eqref{eq:cavityM3comparison}--\eqref{eq:cavityM3endpoint}, this is at
most
\begin{equation}
 C_{\cK}\left(\nu_s[|Q_{12}|^3]+N^{-3/2}\right).
 \label{eq:cavitydiagonalerror}
\end{equation}
Combining \eqref{eq:cavityendpointcomparison},
\eqref{eq:cavitysecondderivativebound}, and
\eqref{eq:cavitydiagonalerror} proves \eqref{eq:rho}.  The coefficient
table \eqref{eq:cavitytable} is therefore equivalent to
\eqref{eq:cavitysystem} with the stated remainder.
\end{proof}

\subsection{Stability of the cavity matrix}\label{subsec:stability}

Put
\begin{equation}
 \gamma'\coloneqq1-4q+3r,\qquad
 \gamma''\coloneqq2q+q^2-3r.
 \label{eq:gammas}
\end{equation}
The change of coordinates
\[
 S\coloneqq
 \begin{pmatrix}
 0&2&-3\\
 1&-4&3\\
 1&-2&1
 \end{pmatrix}
\]
gives, by direct multiplication,
\begin{equation}
 SLS^{-1}
 =
 \begin{pmatrix}
 \beta^2\gamma'&\beta^2\gamma''&0\\
 0&\beta^2\gamma'&0\\
 0&0&\alpha
 \end{pmatrix}.
 \label{eq:triangular}
\end{equation}
Hence
\begin{equation}
 \det(I-sL)
 =(1-s\alpha)(1-s\beta^2\gamma')^2.
 \label{eq:det}
\end{equation}
Since $r\le q$,
\[
 \alpha-\beta^2\gamma'
 =2\beta^2(q-r)\ge0.
\]
Thus $\beta^2\gamma'\le\alpha<1$.  If $\gamma'<0$, then
$1-s\beta^2\gamma'\ge1$; if $\gamma'\ge0$, then
$1-s\beta^2\gamma'\ge1-\alpha$.  Together with
\eqref{eq:delta}, this gives
\begin{equation}
 M_{\cK}\coloneqq
 \sup_{\substack{(\beta,h)\in\cK\\0\le s\le1}}
 \norm{(I-sL)^{-1}}_\infty<\infty.
 \label{eq:inverse}
\end{equation}
Notice that \eqref{eq:triangular}, rather than the eigenvalues alone,
also controls the possible Jordan term.

\section{Conclusions}\label{sec:conclusion}
\subsection{Absorption of the cubic remainder}\label{subsec:absorb}

Recall that $A_s= \nu_s[Q_{12}^2]$. Replica symmetry and Cauchy--Schwarz give
\[
 |B_s|\le A_s,\qquad |C_s|\le A_s.
\]
Therefore $\norm{x_s}_\infty=A_s$.  Inverting
\eqref{eq:cavitysystem}, and using \eqref{eq:inverse} and
\eqref{eq:rho}, yields
\begin{equation}
 A_s
 \le \frac{C_{\cK}}N
 +C_{\cK}N^{-3/2}
 +C_{\cK}\nu_s[|Q_{12}|^3].
 \label{eq:preabsorb}
\end{equation}

Fix $\varepsilon>0$.  Since $|Q_{12}|\le2$,
\begin{align}
 \nu_s[|Q_{12}|^3]
 &\le
 \varepsilon\,\nu_s[Q_{12}^2]
 +8\,\nu_s[|Q_{12}|>\varepsilon]\nonumber\\
 &\le
 \varepsilon A_s
 +8C_{\cK,\varepsilon}e^{-c_{\cK,\varepsilon}N},
 \label{eq:thirdsplit}
\end{align}
where Corollary~\ref{cor:tail} was used.  Choose
$\varepsilon=\varepsilon_{\cK}$
so that the coefficient of $A_s$ obtained by substituting
\eqref{eq:thirdsplit} into \eqref{eq:preabsorb} is at most $1/2$.
We obtain
\[
 A_s\le\frac{C_{\cK}}N
\]
uniformly on $\cK\times[0,1]$, after increasing $C_{\cK}$ to include
the finitely many small values of $N$.  This proves the first assertion
of Theorem~\ref{thm:main}.

\subsection{Free energy}\label{subsec:freeenergy}

The identity \eqref{eq:DNprime} and $D_N(0)=0$ give
\begin{equation}
 \phi^{\mathrm{RS}}(\beta,h)-\phi_N(1)
 =D_N(1)=
 \frac{\beta^2}{4}\int_0^1\nu_s[Q_{12}^2]\,\dd s.
 \label{eq:freeenergyidentity}
\end{equation}
The integrand is nonnegative, and the already proved second-moment
bound makes the right-hand side at most $C_{\cK}/N$.  This proves the
free-energy assertion of Theorem~\ref{thm:main}.

\subsection{Replicon susceptibility}\label{subsec:susceptibility}

Define
\[
 D_s\coloneqq A_s-2B_s+C_s.
\]
The last row of the transformed system \eqref{eq:triangular}, or
equivalently the scalar product of \eqref{eq:cavitysystem} with
$(1,-2,1)$, gives
\begin{equation}
 (1-s\alpha)D_s=\frac{\alpha}{\beta^2N}+\rho^{(R)}_{N,s},
 \label{eq:repliconeq}
\end{equation}
where
\begin{equation}
 |\rho^{(R)}_{N,s}|
 \le C_{\cK}\left(N^{-3/2}+\nu_s[|Q_{12}|^3]\right).
 \label{eq:repliconrho}
\end{equation}
Here
\[
 (1,-2,1)\theta=1-2q+r=\frac{\alpha}{\beta^2}.
\]

The $O(N^{-1})$ second-moment estimate and
\eqref{eq:thirdsplit} imply, for every fixed $\varepsilon>0$,
\[
 \sup_{\cK\times[0,1]}
 N\nu_s[|Q_{12}|^3]
 \le C_{\cK}\varepsilon
 +8N C_{\cK,\varepsilon}e^{-c_{\cK,\varepsilon}N}.
\]
First let $N\to\infty$ and then $\varepsilon\downarrow0$.  Thus
\[
 \nu_s[|Q_{12}|^3]=o_{\cK}(N^{-1})
\]
uniformly along the path.  It follows from
\eqref{eq:repliconeq}, \eqref{eq:repliconrho}, and
$1-s\alpha\ge\delta_{\cK}$ that
\[
 ND_s
 =
 \frac{\alpha}{\beta^2(1-s\alpha)}+o_{\cK}(1),
\]
which is the final assertion of Theorem~\ref{thm:main}.

\appendix

\section{The specialized Guerra--Talagrand bound}
\label{app:specialGT}

The general two-dimensional Guerra--Talagrand bound is classical; see,
in particular, \cite{ChenGT} and \cite{Talagrand}.  For completeness, we
verify here the precise finite-step specialization and normalization
used in the proof of Proposition~\ref{prop:GTcoercivity}.  Retaining the
proof records the rank-one covariance pieces, the factors $1/2$, and
the Lagrange-multiplier convention for the present path.  No claim of
novelty is made for the abstract Guerra--Talagrand interpolation theorem.

We first record the Gaussian covariance interpolation identity used in
the verification.  If $G^{(0)}$ and $G^{(1)}$ are independent centered
Gaussian vectors with covariance matrices $C_0$ and $C_1$, respectively,
and $G_t\coloneqq\sqrt t\,G^{(1)}+\sqrt{1-t}\,G^{(0)}$, then applying
\eqref{eq:GIBP} twice yields, for $0<t<1$,
\begin{equation}
 \frac{\dd}{\dd t}\E F(G_t)
 =\frac12\sum_{i,j}(C_1-C_0)_{ij}
   \E[\partial_{ij}F(G_t)].
 \label{eq:gaussianinterpolation}
\end{equation}
The right-hand side has a continuous extension to $t=0,1$.

\begin{lemma}[GT bound for the replica-symmetric two-replica path]
\label{lem:specialGT}
If the centered one-replica Hamiltonian has covariance
$N\xi_s(R(\sigma,\tau))$, then, for every $v\in\mathcal R_N$ and
$\lambda\in\mathbb R$,
\begin{equation}
 \frac1N\E\log Z^{(2)}_{N,s}(v)
 \le \mathfrak P_s(\lambda,v).
 \label{eq:specialGTlemma}
\end{equation}
\end{lemma}

\begin{proof}
We use only the three Gaussian recursion operators
\[
 T_0F\coloneqq\E F,\qquad
 T_{1/2}F\coloneqq2\log\E e^{F/2},\qquad
 T_1F\coloneqq\log\E e^F.
\]
Let $0=u_0<\cdots<u_K=1$ be the partition obtained from the distinct
points among $0,q,|v|,1$.  Thus $K\le3$, and the value $m_k$ of
$\gamma_v$ on $[u_{k-1},u_k)$ belongs to
$\{0,\frac12,1\}$.  Put, entrywise,
\[
 C_k\coloneqq\xi_s'(Q^v_{u_k}),\qquad
 \Theta(x)\coloneqq x\xi_s'(x)-\xi_s(x),\qquad
 D_k\coloneqq\sum_{a,b=1}^2\Theta(Q^v_{u_k,ab}).
\]
By \eqref{eq:GTspeed},
\begin{align*}
 C_k-C_{k-1}&=\int_{u_{k-1}}^{u_k}A_s^v(u)\,\dd u\succeq0,\\
 D_k-D_{k-1}&=\int_{u_{k-1}}^{u_k}
 \sum_{a,b=1}^2A^v_{s,ab}(u)Q^v_{u,ab}\,\dd u\ge0.
\end{align*}
Take independent Gaussian vector increments $z_{i,k}$ with covariance
$C_k-C_{k-1}$ and scalar increments $y_k$ with variance
$D_k-D_{k-1}$.  At each site the base vector has covariance
$C_0=\beta^2(1-s)q\,\mathbf1\mathbf1^{\mathsf T}$ and hence equals
$(Y_{i,0}-h,Y_{i,0}-h)$, where the $Y_{i,0}$ are independent copies
of $Y_0$.

For a constrained pair
$\boldsymbol\sigma=(\sigma^1,\sigma^2)$ and $0\le t\le1$, set
\begin{align*}
 \mathcal H_t(\boldsymbol\sigma)
 \coloneqq{}&\sqrt t\,[H_{N,s}^{\rm cen}(\sigma^1)
                     +H_{N,s}^{\rm cen}(\sigma^2)]\\
 &+\sqrt{1-t}\sum_{i=1}^N\sum_{a=1}^2
 \left(Y_{i,0}-h+\sum_{k=1}^Kz_{i,k}^a\right)\sigma_i^a
 +\sqrt{tN}\sum_{k=1}^K y_k\\
 &+h\sum_{i=1}^N(\sigma_i^1+\sigma_i^2)
 +\lambda\sum_{i=1}^N\sigma_i^1\sigma_i^2-\lambda Nv.
\end{align*}
Let $X_K(t)$ be the logarithm of the sum of
$e^{\mathcal H_t(\boldsymbol\sigma)}$ over $R_{12}=v$.  Integrate the
increments backward by
\[
 X_{k-1}(t)=T_{m_k}X_k(t),\qquad k=K,\ldots,1,
\]
and set $\varphi_N(t)=N^{-1}\E X_0(t)$, where the last expectation is
over the original disorder and the base fields.  This is a finite
recursion, including when $m_k=0$ or $1$.

At $t=0$, remove the overlap constraint.  The recursion then
factorizes over sites, and the Cole--Hopf formula on each interval
gives
\begin{equation}
 \varphi_N(0)\le2\log2+
 \E U_s^{\lambda,v}(0,Y_0,Y_0)-\lambda v.
 \label{eq:GTendpointzero}
\end{equation}
At $t=1$, the spin Hamiltonian is independent of the auxiliary
increments, while
\[
 T_{m_k}(x+\sqrt N\,y_k)
 =x+\frac{Nm_k}{2}(D_k-D_{k-1}).
\]
Therefore
\begin{equation}
 \varphi_N(1)=\frac1N\E\log Z^{(2)}_{N,s}(v)+\mathcal L_s(v).
 \label{eq:GTendpointone}
\end{equation}

For completeness, repeated Gaussian integration by parts through the
finite recursion gives its monotonicity.  Set $m_0=0$, $m_{K+1}=1$.
The chain rule writes $\varphi_N'(t)$ as
\begin{equation}
 -\frac12\sum_{k=0}^{K}(m_{k+1}-m_k)
 \E\left\langle
 \sum_{a,b=1}^2
 \Delta_{\xi_s}(R^{ab},Q^v_{u_k,ab})
 \right\rangle_{k,t},
 \label{eq:GTderivative}
\end{equation}
where $\langle\cdot\rangle_{k,t}$ is the two-copy Gibbs expectation
produced by differentiating the $k$th nested expectation.  This
formula follows directly from \eqref{eq:gaussianinterpolation}; the
covariance derivative for each pair $(a,b)$ is
\[
 \Delta_{\xi_s}(x,y)
 \coloneqq\xi_s(x)-\xi_s(y)-\xi_s'(y)(x-y)
 =\frac{s\beta^2}{2}(x-y)^2\ge0.
\]
The coefficients in \eqref{eq:GTderivative} are nonnegative because
$\gamma_v$ is nondecreasing.  Hence $\varphi_N(1)\le\varphi_N(0)$.
Combining the two endpoint identities and using
\eqref{eq:GTfunctional} proves \eqref{eq:specialGTlemma}.
\end{proof}

\section{Lean 4 formalization guide}
\label{app:lean-guide}

This section is a working specification for a kernel-checked Lean 4
formalization of the paper.  It is deliberately more explicit than the
mathematical proof.  It records proposed data types, theorem interfaces,
normal forms, dependency order, finite-size conventions, and the analytic
facts that must be supplied.  The intended target is a project using a
fixed release of Lean 4 and Mathlib.  Every final declaration should be
free of \texttt{sorry}; no axiom is needed or intended.

Code fragments below are design templates.  Names for Gaussian measures,
Bochner integrals, and differentiability lemmas can change between Mathlib
releases.  Before copying a fragment, inspect the selected release with
\texttt{\#check}.  Definitions local to this project should keep the names
used below even if the imported API differs.  That thin wrapper will make
the rest of the development stable under library updates.

\subsection{Exact formalization target and trusted boundary}

The strongest target is a theorem whose assumptions contain only the
definitions appearing in the paper: the SK disorder consists of independent
standard Gaussian coordinates, $q$ solves \eqref{eq:q}, $r$ and $\alpha$ are
given by \eqref{eq:ra}--\eqref{eq:alpha}, and the parameter set is a compact
subset of the strict AT region.  Its conclusion is the three-part statement
of Theorem~\ref{thm:main}, with all uniform quantifiers displayed.

It is useful to maintain two theorem layers.
The pointwise layer fixes $(\beta,h,q,r,\alpha)$ and explicit positive
constants witnessing every estimate.  The uniform layer derives common
constants from compactness and continuity.  Nearly all probability and
matrix calculations belong to the pointwise layer.  Compactness should be
introduced only after those calculations have been checked.

The formalization may import proved Mathlib results about finite sums,
matrices, derivatives, integrals, Gaussian random variables, conditional
expectation, concentration, and Gronwall inequalities.  If a required
result is absent, add it as a proved project lemma.  Do not encode a standard
tool as a field of the final theorem merely to avoid proving it.  Temporary
interface hypotheses are acceptable during development, but the final
dependency report should show how each one is discharged.

The following statements are genuine mathematical inputs rather than
mere simplifications:
\begin{itemize}[leftmargin=2em]
\item existence and uniqueness of the positive-field fixed point $q$;
\item Gaussian integration by parts and covariance interpolation;
\item Gaussian concentration for a finite maximum;
\item It\^o's formula and the relevant backward Kolmogorov formula;
\item the finite-recursion Guerra--Talagrand interpolation inequality;
\item differentiation under the finite-volume disorder integral;
\item the finite-dimensional Gronwall inequalities used in
      Sections~\ref{subsec:gronwall} and \ref{subsec:cavity}.
\end{itemize}
Each item must end as an imported theorem or a proved local theorem.  The
paper's label 'standard' is not itself a formal justification.

\subsection{Dependency graph and recommended order}

The proof dependency graph is
\[
\begin{array}{c}
\text{finite spins, finite Gibbs expectations, Gaussian wrappers}\\[-0.2em]
\Downarrow\\[-0.2em]
\text{fixed-point moments and elementary hyperbolic identities}
\end{array}
\]
\[
\begin{array}{c}
\text{scalar semigroup and local-field identities}\\[-0.2em]
\Downarrow\\[-0.2em]
\text{strict AT sign}\\[-0.2em]
\Downarrow\\[-0.2em]
\text{two-replica GT coercivity}\\[-0.2em]
\Downarrow\\[-0.2em]
\text{preliminary concentration and fixed tails},
\end{array}
\]
followed by
\[
\begin{array}{c}
\text{cavity derivative and coefficient table}\\[-0.2em]
\Downarrow\\[-0.2em]
\text{matrix stability}\\[-0.2em]
\Downarrow\\[-0.2em]
\text{cubic absorption}\\[-0.2em]
\Downarrow\\[-0.2em]
\text{main second-moment bound},
\end{array}
\]
and finally
\[
\begin{array}{c}
\text{smart-path identity and replicon row}\\[-0.2em]
\Downarrow\\[-0.2em]
\text{free-energy and susceptibility conclusions}.
\end{array}
\]

The most productive initial milestone is a completely checked
finite-algebra package containing spin sums, overlap bounds, the cavity
coefficient table, the triangularization of $L$, and the final absorption
lemma.  These results expose normalization errors early and do not depend on
stochastic calculus.

\subsection{Proposed file layout}

A practical project layout is as follows.  The comments state the intended
contents, not extra assumptions.
\begin{lstlisting}
SKKusuoka/
  Basic.lean              -- scalar inequalities and project conventions
  Gaussian.lean           -- Gaussian expectation and IBP wrappers
  Spins.lean              -- Ising configurations and replica overlaps
  Gibbs.lean              -- finite partition functions and Gibbs averages
  Parameters.lean         -- q, r, alpha, strict AT data
  SmartPath.lean          -- H_{N,s}, phi_N, endpoint and covariance facts
  ScalarPDE.lean          -- explicit finite-step scalar semigroup
  Diffusion.lean          -- local field, Ito identities, AT sign
  Matrix3.lean            -- theta, L, S, inverse bounds
  GT/Terminal.lean        -- f_lambda and derivative bounds
  GT/FiniteRecursion.lean -- two-dimensional finite semigroup
  GT/Interpolation.lean   -- specialized GT inequality
  GT/Coercivity.lean      -- quadratic constrained-pressure loss
  Concentration.lean      -- coupled pressure, Gronwall, fixed tails
  Cavity/Operator.lean    -- replica operator D_n
  Cavity/Coefficients.lean-- exact coefficient enumeration
  Cavity/Expansion.lean   -- Taylor remainder and cavity system
  Absorption.lean         -- cubic closure
  Main.lean               -- paper-level theorem and consequences
\end{lstlisting}

Keep imports acyclic.  In particular, \texttt{Matrix3.lean} should import
only elementary parameter facts, and \texttt{Cavity/Coefficients.lean}
should not import any concentration theorem.  The coefficient calculation
is a finite identity and should remain testable in isolation.

\subsection{Naming, normalization, and simp policy}

Use ASCII declaration names even when the mathematical notation is Greek.
Suggested translations include
\begin{center}
\begin{tabular}{lll}
\toprule
paper notation & Lean name & type or role\\
\midrule
$\beta,h,q,r,\alpha$ & \texttt{beta, h, q, r, alpha} & real parameters\\
$\Sigma_N$ & \texttt{SpinConfig N} & finite configuration type\\
$R_{ab}$ & \texttt{overlap N sigma a b} & uncentered overlap\\
$Q_{ab}$ & \texttt{centeredOverlap p sigma a b} & $R_{ab}-q$\\
$\langle f\rangle_s$ & \texttt{gibbsExpect path f} & finite Gibbs average\\
$\nu_s[f]$ & \texttt{nu path f} & disorder-averaged Gibbs value\\
$P_s^*$ & \texttt{rsPathValue p s} & scalar trial value\\
$A_s,B_s,C_s$ & \texttt{momentA, momentB, momentC} & overlap moments\\
$\mathcal D_n$ & \texttt{cavityOp n} & replica derivative operator\\
\bottomrule
\end{tabular}
\end{center}

Choose one normalization for expectation and never switch silently between
a probability expectation and an unnormalized integral.  Likewise, define
the Hamiltonian with the factor $N^{-1/2}$ exactly once.  A separate lemma
should state the centered covariance \eqref{eq:GTcovariance}, including the
finite constant $-s\beta^2/2$.

Tag only stable canonical facts with \texttt{[simp]}.  Good candidates are
the square of a spin, overlap symmetry, the diagonal overlap, evaluation of
the cavity endpoint, and obvious matrix-entry reductions.  Do not mark
fixed-point rewrites in both directions.  Prefer the rewrite orientation
$\E\tanh^2(h+\beta\sqrt q Z)\mapsto q$, since this reduces analytic
expressions to parameters.

\subsection{Finite spin space and overlaps}

Using \texttt{Bool} for a spin makes the configuration type finite without
having to build a new two-point \texttt{Fintype}.  Convert a Boolean to a
real spin only at the algebraic boundary.
\begin{lstlisting}
import Mathlib

open scoped BigOperators

namespace SK

def spinVal : Bool -> Real
  | false => -1
  | true  => 1

@[simp] theorem spinVal_sq (b : Bool) : spinVal b ^ 2 = 1 := by
  cases b <;> norm_num [spinVal]

abbrev SpinConfig (N : Nat) := Fin N -> Bool
abbrev Replicas (N k : Nat) := Fin k -> SpinConfig N

def overlap (N : Nat) (sigma tau : SpinConfig N) : Real :=
  (N : Real)⁻¹ * Finset.univ.sum
    (fun i : Fin N => spinVal (sigma i) * spinVal (tau i))

def replicaOverlap (N : Nat) (sigmas : Replicas N k)
    (a b : Fin k) : Real := overlap N (sigmas a) (sigmas b)

def centeredOverlap (q : Real) (N : Nat) (sigmas : Replicas N k)
    (a b : Fin k) : Real := replicaOverlap N sigmas a b - q

end SK
\end{lstlisting}

The paper is stated for $N\ge1$.  Carry a hypothesis \texttt{hN : 0 < N}
in every lemma that cancels $N$ or uses the diagonal identity.  This is
safer than placing positivity inside the definition of
\texttt{SpinConfig}.  Prove the following facts immediately:
\begin{align*}
R(\sigma,\tau)&=R(\tau,\sigma),&
R(\sigma,\sigma)&=1,\\
|R(\sigma,\tau)|&\le1,&
|Q(\sigma,\tau)|&\le2
\quad\text{when }0\le q\le1.
\end{align*}
Also prove the exact range characterization
\[
 R(\sigma,\tau)=1-\frac{2k}{N}
 \quad\text{for some }k\in\{0,\ldots,N\}.
\]
This identifies $\mathcal R_N$, proves $|\mathcal R_N|\le N+1$, and
supports the finite maximum in \eqref{eq:gaussianmax}.  It is preferable to
define the range as an image rather than as a set described by an
existential:
\begin{lstlisting}
def overlapRange (N : Nat) : Finset Real :=
  Finset.univ.image (fun k : Fin (N + 1) =>
    1 - 2 * (k : Real) / N)
\end{lstlisting}
The equality between this finite set and the actual overlap image needs the
Hamming-distance counting lemma.  For $N=0$, keep a harmless formal
definition but never invoke the range theorem.

\subsection{Finite Gibbs measures without measure-theoretic overhead}

For fixed disorder, all Gibbs operations are finite sums.  Define them as
such; do not introduce a measure on the finite spin space unless a later API
requires it.
\begin{lstlisting}
namespace SK

noncomputable def partition (H : SpinConfig N -> Real) : Real :=
  Finset.univ.sum (fun sigma => Real.exp (H sigma))

noncomputable def gibbsWeight (H : SpinConfig N -> Real)
    (sigma : SpinConfig N) : Real :=
  Real.exp (H sigma) / partition H

noncomputable def gibbsExpect (H : SpinConfig N -> Real)
    (f : SpinConfig N -> Real) : Real :=
  Finset.univ.sum (fun sigma => gibbsWeight H sigma * f sigma)

noncomputable def replicaGibbsExpect (H : SpinConfig N -> Real)
    (f : Replicas N k -> Real) : Real :=
  Finset.univ.sum (fun sigmas =>
    (Finset.univ.prod (fun a : Fin k => gibbsWeight H (sigmas a))) *
      f sigmas)

end SK
\end{lstlisting}

The foundational lemmas should say
\begin{itemize}[leftmargin=2em]
\item \texttt{0 < partition H};
\item Gibbs weights are nonnegative and sum to one;
\item Gibbs expectation is linear, monotone, and preserves constants;
\item the replica expectation factors for observables depending on disjoint
      replica coordinates;
\item permutation of replica labels leaves the expectation unchanged;
\item if $|f|\le C$ pointwise, then
      $|\langle f\rangle|\le C$.
\end{itemize}
The positivity proof is short because every exponential is positive and the
configuration type is nonempty.  State a reusable lemma for positive finite
sums so later denominator arguments close by \texttt{positivity} or
\texttt{field\_simp}.

For constrained sums, filter rather than use a subtype:
\begin{lstlisting}
noncomputable def constrainedPartition2
    (H : SpinConfig N -> Real) (v : Real) : Real :=
  Finset.univ.sum (fun st : SpinConfig N × SpinConfig N =>
    if overlap N st.1 st.2 = v then
      Real.exp (H st.1 + H st.2)
    else 0)
\end{lstlisting}
This convention makes an unavailable overlap value give partition function
zero.  The logarithm is used only under the hypothesis
\texttt{v in overlapRange N}; prove positivity under that hypothesis before
forming \texttt{Real.log}.

\subsection{Parameter records and the fixed-point layer}

Do not ask Lean to infer every parameter restriction from membership in a
large set.  A record gives convenient named hypotheses for the pointwise
proof.
\begin{lstlisting}
structure RSPoint where
  beta : Real
  h : Real
  q : Real
  r : Real
  beta_pos : 0 < beta
  h_pos : 0 < h
  q_nonneg : 0 <= q
  q_lt_one : q < 1
  r_nonneg : 0 <= r
  r_le_q : r <= q
  fixedPoint : q = gaussianTanhSq beta h q
  fourthMoment : r = gaussianTanhFourth beta h q

namespace RSPoint

def alpha (p : RSPoint) : Real :=
  p.beta ^ 2 * (1 - 2 * p.q + p.r)

def StrictAT (p : RSPoint) : Prop := p.alpha < 1

end RSPoint
\end{lstlisting}
Here \texttt{gaussianTanhSq} and \texttt{gaussianTanhFourth} are project
wrappers around integration against a standard normal law.  The record does
not create new assumptions: a constructor theorem must build it from
$(\beta,h)$ and the proved existence and uniqueness theorem for $q$.

A second record can package a uniform gap when working on a fixed compact
set:
\begin{lstlisting}
structure UniformATData where
  K : Set (Real × Real)
  compact_K : IsCompact K
  beta_pos : forall x, x ∈ K -> 0 < x.1
  h_pos : forall x, x ∈ K -> 0 < x.2
  delta : Real
  delta_pos : 0 < delta
  alpha_gap : forall x, x ∈ K -> alphaAt x <= 1 - delta
\end{lstlisting}
The final uniform theorem should construct \texttt{delta} from compactness
and continuity instead of requiring the user to supply it.  Keeping this
record available is still useful for the long middle part of the proof.

\subsection{Gaussian expectation wrappers and analytic contracts}

Create one project definition for the standard Gaussian law and one for
expectation.  This confines any Mathlib API adjustment to a small file.
The desired semantic equations are
\[
 \operatorname{gaussE}(f)=\int_{\R} f(x)\,\gamma(\dd x),
 \qquad \gamma=N(0,1),
\]
and
\[
 \operatorname{gaussianTanhSq}(\beta,h,q)
 =\operatorname{gaussE}\bigl(x\mapsto
   \tanh^2(h+\beta\sqrt q\,x)\bigr).
\]
The wrapper file should export these facts:
\begin{itemize}[leftmargin=2em]
\item the Gaussian law is a probability measure;
\item bounded continuous functions are integrable;
\item expectation is linear and monotone;
\item dominated convergence in a parameterized form;
\item the standard normal has mean zero and variance one;
\item a sum of independent centered Gaussians is Gaussian, with variances
      adding;
\item the integration-by-parts identity \eqref{eq:GIBP};
\item covariance differentiation in the form
      \eqref{eq:gaussianinterpolation}.
\end{itemize}

State integration by parts first for a finite index type.  This matches all
finite-volume applications and avoids committing to Euclidean coordinate
notation.
\begin{lstlisting}
theorem gaussian_ibp_fin
    {m : Nat} (G : Omega -> (Fin m -> Real))
    (C : Matrix (Fin m) (Fin m) Real)
    (hG : CenteredGaussianVector G C)
    (F : (Fin m -> Real) -> Real)
    (hF : ContDiff Real 1 F)
    (hInt : GaussianIBPIntegrable G F) (i : Fin m) :
    E (fun w => G w i * F (G w)) =
      Finset.univ.sum (fun j =>
        C i j * E (fun w =>
          fderiv Real F (G w) (Pi.single j 1))) := by
  -- Apply the selected library theorem, or prove density IBP here.
  ...
\end{lstlisting}
The three dots mark a proof obligation, not a permitted final term.
In actual source, replace them before any downstream file is declared
complete.  It is worth defining a project predicate such as
\texttt{GaussianIBPIntegrable} to gather all integrability conditions used by
the library theorem.  Prove small automation lemmas showing that finite sums
of exponentials of affine Gaussian functions satisfy it.

\subsection{Hyperbolic identities and moment inequalities}

Before formalizing the fixed-point theorem, build a small deterministic
library for \texttt{Real.tanh} and \texttt{Real.cosh}.  Define
\[
 \operatorname{sech}(x)=\frac1{\cosh x}
\]
locally if Mathlib has no stable project-facing name.  Required identities
include
\begin{align*}
&\cosh x>0,\qquad |\tanh x|<1,\qquad 0<\operatorname{sech}x,\\
&1-\tanh^2x=\operatorname{sech}^2x,\\
&(1-\tanh^2x)^2
 =1-2\tanh^2x+\tanh^4x.
\end{align*}
From $|\tanh x|\le1$, prove $\tanh^4x\le\tanh^2x$.  Integration then gives
$0\le r\le q\le1$.  The identity in \eqref{eq:alpha} should be a named
lemma:
\begin{lstlisting}
theorem alpha_eq_gaussian_sech_fourth (p : RSPoint) :
    p.alpha =
      p.beta ^ 2 * gaussianSechFourth p.beta p.h p.q := by
  rw [RSPoint.alpha, p.fixedPoint, p.fourthMoment]
  -- Push the pointwise polynomial identity through the integral.
  ...
\end{lstlisting}
Also isolate the strictly positive conclusions $0<q$ and $q<1$ for $h>0$.
The proof of $q<1$ uses strict almost-everywhere inequality plus probability
mass one; plain monotonicity alone is insufficient.

For existence and uniqueness of $q$, define
\[
 T_{\beta,h}(u)=\E\tanh^2(h+\beta\sqrt u\,Z),
 \qquad 0\le u\le1.
\]
Separate the following lemmas: continuity of $T$, preservation of $[0,1]$,
existence of a fixed point, positive-field uniqueness, and continuous
dependence of the unique fixed point on $(\beta,h)$.  This separation is
important because only continuity and uniqueness are needed for the compact
uniformity argument.  If the uniqueness proof is imported from another
formal development, state its exact hypotheses and verify that $h>0$ matches
them.

\subsection{Smart-path data and endpoint checks}

Represent all disorder coordinates on one probability space.  A convenient
abstract interface is a structure containing the pair couplings, site
fields, independence, and standard-normal proofs.  Later instantiate it by a
finite product Gaussian space.
\begin{lstlisting}
structure SKDisorder (N : Nat) (Omega : Type*)
    [MeasurableSpace Omega] (mu : Measure Omega) where
  g : Fin N -> Fin N -> Omega -> Real
  z : Fin N -> Omega -> Real
  g_symm : forall i j w, g i j w = g j i w
  g_diag : forall i w, g i i w = 0
  measurable_g : forall i j, Measurable (g i j)
  measurable_z : forall i, Measurable (z i)
  gaussian_family :
    JointStandardGaussianUpperTriangleAndSites g z mu
\end{lstlisting}
The final field should itself be defined in \texttt{Gaussian.lean}; it must
express the full covariance matrix, which encodes both independence and unit
variance.  An upper-triangular index type is cleaner than storing a symmetric
array, but either design is acceptable if the covariance calculation stays
simple.

Define the path Hamiltonian directly from \eqref{eq:path}.  Name and prove:
\begin{itemize}[leftmargin=2em]
\item measurability and integrability of $H_{N,s}(\sigma)$;
\item positivity and integrability of its partition function;
\item the $s=0$ factorization into independent one-site terms;
\item the $s=1$ equality with the original Hamiltonian;
\item the centered covariance formula \eqref{eq:GTcovariance};
\item continuity in $s$ of every finite-volume expectation later used;
\item $D_N(0)=0$ after identifying the $s=0$ free energy with $P_0^*$.
\end{itemize}

Avoid differentiating square roots at $s=0$ or $s=1$.  Prove the derivative
identity first on the open interval and extend it to the endpoints by the
finite-dimensional dominated-convergence lemma, exactly as in the paper.
This pattern is also required for the cavity parameter $u$.

\subsection{Replica indices and exchangeability}

The cavity proof repeatedly appends replicas.  Fixed types such as
\texttt{Fin 2}, \texttt{Fin 3}, and \texttt{Fin 4} work for final moments,
but the derivative operator needs a systematic embedding
$\operatorname{Fin} n\hookrightarrow\operatorname{Fin}(n+2)$.
Use Mathlib's canonical \texttt{Fin.castLE}, \texttt{Fin.castSucc}, and
\texttt{Fin.last} maps or define project wrappers with simp lemmas.

Prove exchangeability once for arbitrary permutations:
\begin{lstlisting}
theorem replica_expect_perm
    (pi : Equiv.Perm (Fin k)) (f : Replicas N k -> Real) :
    replicaGibbsExpect H (fun sigmas =>
      f (fun a => sigmas (pi a))) = replicaGibbsExpect H f := by
  -- Reindex by the induced equivalence on replica tuples.
  ...
\end{lstlisting}
Pairwise moment identifications, H\"older estimates, and coefficient-table
symmetries should then be consequences of this theorem, not fresh finite-sum
proofs.

\subsection{Scalar PDE as an explicit semigroup}

The paper already supplies the solution
\eqref{eq:PDEupper}--\eqref{eq:PDElower}.  For formalization, define $\Psi$
by those formulas and prove that it solves the PDE.  This is easier to audit
than starting from an abstract PDE existence theorem.

Define two branches:
\begin{align*}
 \Psi_{\rm up}(u,x)
 &=\log\cosh x+\frac{s\beta^2}{2}(1-u),\\
 \Psi_{\rm low}(u,x)
 &=\E\log\cosh\bigl(x+\beta\sqrt{s(q-u)}Z\bigr)
   +\frac{s\beta^2}{2}(1-q).
\end{align*}
Then define \texttt{psi p s u x} by a decidable branch on $q\le u$.
Prove branch agreement at $u=q$, continuity there, and the derivatives needed
on each open subinterval.  The target regularity is piecewise $C^{1,2}$;
global twice differentiability in $u$ is neither asserted nor needed.

The following derivative formulas deserve individual names:
\begin{align*}
 \Psi_x(u,x)&=\tanh x &&(q\le u),\\
 \Psi_{xx}(u,x)&=\operatorname{sech}^2x &&(q\le u),\\
 \Psi_x(u,x)&=\E\tanh(x+\beta\sqrt{s(q-u)}Z)&&(u<q),\\
 \Psi_{xx}(u,x)&=\E\operatorname{sech}^2
     (x+\beta\sqrt{s(q-u)}Z)&&(u<q).
\end{align*}
Record uniform bounds $|\Psi_x|\le1$ and $0\le\Psi_{xx}\le1$.  They discharge
most integrability and Lipschitz obligations for the local-field SDE.

The reduction of \eqref{eq:scalartrial} to \eqref{eq:RSpathvalue} should be
a separate algebraic theorem.  Its critical variance identity is
\[
 \beta^2(1-s)q+s\beta^2q=\beta^2q.
\]
Normalize polynomial expressions with \texttt{ring}; do not ask
\texttt{simp} to find this cancellation.

\subsection{Local field, diffusion, and strict AT sign}

The formal local field is an SDE with bounded Lipschitz drift
\[
 b(u,x)=s\beta^2x_q(u)\Psi_x(u,x)
\]
and constant diffusion coefficient $\beta\sqrt s$.  Because $x_q$ has one
jump, either invoke an SDE theorem for measurable time-dependent drift or
construct the process in two pieces and concatenate at $q$.  The second
route aligns with the explicit semigroup and usually requires fewer general
measurability results.

The essential distribution statement is \eqref{eq:Xq}.  Prove it as a
named theorem from the sum-of-independent-Gaussians lemma.  Downstream
moment rewrites should use this theorem rather than reopening the SDE
construction.

For $u<q$, prove
\[
 g_s'(u)=s\beta^2\E\bigl[\Psi_{xx}(u,X_u)^2\bigr]
\]
and use $0\le\Psi_{xx}\le1$ plus the appropriate comparison at $q$ to obtain
\eqref{eq:leftstrict}.  For $u\ge q$, define $m_u=\tanh X_u$ and apply
It\^o's formula.  The drift cancellation should end in the exact normal form
\[
 \dd m_u=\beta\sqrt s(1-m_u^2)\dd W_u.
\]
If Lean leaves a residual expression, first rewrite
$\tanh'(x)=1-\tanh^2x$ and
$\tanh''(x)=-2\tanh x(1-\tanh^2x)$, then use \texttt{ring}.

The diffusion comparison Lemma~\ref{lem:diffusioncomparison} should be split
into the following reusable results:
\begin{itemize}[leftmargin=2em]
\item a time-change representation of the martingale $m$;
\item the representation \eqref{eq:Flambda-representation};
\item positivity and unit mean of the density $F_\lambda$;
\item its monotonicity property needed for the comparison;
\item the moment inequality
      $\beta^2\E(1-m_u^2)^2\le\alpha$.
\end{itemize}
The assumption $s\beta^2(1-q)>1$ is used only in this comparison branch.
Keep it local so the complementary elementary branch remains easy to apply.

Package Proposition~\ref{prop:pathRS} into sign, local slope, and compact
gap theorems.  A useful interface is
\begin{lstlisting}
theorem gs_sub_id_sign (p : RSPoint) (hat : p.StrictAT)
    (hs : s ∈ Set.Icc (0 : Real) 1) :
    (forall u ∈ Set.Ico (0 : Real) p.q, 0 < gs p s u - u) ∧
    gs p s p.q - p.q = 0 ∧
    (forall u ∈ Set.Ioc p.q 1, gs p s u - u < 0) := by
  ...

theorem gs_linear_gap_on_compact
    (U : UniformATData) :
    exists c > 0, exists eps > 0,
      forall x ∈ U.K, forall s ∈ Set.Icc (0 : Real) 1,
      forall u ∈ Set.Icc (0 : Real) 1,
      abs (u - qAt x) <= eps ->
      abs (gsAt x s u - u) >= c * abs (u - qAt x) := by
  ...
\end{lstlisting}
In the actual file, eliminate every placeholder.  The second theorem should
be derived from the pointwise slope at $q$, continuity of the derivative,
the common AT gap, and a finite subcover or compact minimum argument.

\subsection{Two-replica terminal function}

Define the terminal function without expanding the four-state sum:
\begin{lstlisting}
noncomputable def gtTerminal (lambda x1 x2 : Real) : Real :=
  Real.log ((1 / 4 : Real) *
    ((Finset.univ : Finset Bool).sum (fun e1 =>
      (Finset.univ : Finset Bool).sum (fun e2 =>
        Real.exp (spinVal e1 * x1 + spinVal e2 * x2 +
          lambda * spinVal e1 * spinVal e2)))))
\end{lstlisting}
Prove the inside of the logarithm is positive.  Then derive, rather than
postulate, the useful closed form
\[
 f_\lambda(x_1,x_2)
 =\log\bigl(
 \cosh x_1\cosh x_2\cosh\lambda
 +\sinh x_1\sinh x_2\sinh\lambda\bigr).
\]
The finite-sum definition is best for spin differentiation and uniform
bounds; the closed form is best for continuity.

The terminal identities required later are
\[
 f_0(x_1,x_2)=\log\cosh x_1+\log\cosh x_2,
 \qquad
 \partial_\lambda f_\lambda\big|_{\lambda=0}
 =\tanh x_1\tanh x_2.
\]
Also prove
\[
 |\partial_\lambda f_\lambda|\le1,\qquad
 0\le\partial_{\lambda\lambda}f_\lambda\le1.
\]
The second derivative is the variance of the two-spin product under a
four-state Gibbs law.  Formalizing it in that form often avoids a large
quotient derivative.

\subsection{Matrix-valued GT path}

Use \texttt{Matrix (Fin 2) (Fin 2) Real} for the path $Q^v$ and speed
$A_s^v$.  Define the sign $\iota_v$ separately, and decide once what happens
at $v=0$.  The choice is irrelevant because the off-diagonal value is zero,
but a simp theorem should state this fact.

The path has a breakpoint $r_v=|v|$.  Required elementary facts include
\[
0\le r_v\le1,\qquad
Q^v_0=0,\qquad
Q^v_1=\begin{pmatrix}1&v\\v&1\end{pmatrix}.
\]
Prove positive semidefiniteness by completing the square:
\[
 (x,y)
 \begin{pmatrix}u&\iota_vu\\\iota_vu&u\end{pmatrix}
 \binom{x}{y}
 =u(x+\iota_vy)^2
\]
below the breakpoint, and
\[
 u(x^2+y^2)+2vxy
 =\frac{u+v}{2}(x+y)^2+\frac{u-v}{2}(x-y)^2
\]
above it.  The hypotheses $u\ge|v|$ make both coefficients nonnegative.
These scalar identities are often easier for Lean than invoking an
eigenvalue criterion.

For the correction term, expand the finite sum over \texttt{Fin 2} and prove
\[
 \mathcal L_s(v)=s\beta^2\int_0^1x_q(u)u\,\dd u
 =\frac{s\beta^2}{2}(1-q^2).
\]
The proof must retain the factor $1/2$ in $\gamma_v$ below $r_v$.  Give the
final independence of $v$ its own theorem:
\begin{lstlisting}
theorem gtCorrection_eq (p : RSPoint)
    (hs : s ∈ Set.Icc (0 : Real) 1)
    (hv : v ∈ Set.Icc (-1 : Real) 1) :
    gtCorrection p s v = s * p.beta ^ 2 / 2 * (1 - p.q ^ 2) := by
  ...
\end{lstlisting}

\subsection{Finite-semigroup construction of the two-dimensional PDE}

The path coefficients are piecewise constant or affine with finitely many
breakpoints drawn from $\{0,q,|v|,1\}$.  Sort and deduplicate this list.
Do not define a recursion indexed by a list whose order has not been proved.

For a positive semidefinite covariance matrix $D$ and $m\ge0$, use the
semigroup
\[
 T_{m,D}F(x)=
 \begin{cases}
 m^{-1}\log\E\exp(mF(x+G_D)),&m>0,\\
 \E F(x+G_D),&m=0.
 \end{cases}
\]
Create a structure for a finite path containing the ordered breakpoints,
increment covariances, nonnegative masses, and proofs of positive
semidefiniteness.  The recursion defining $U_s^{\lambda,v}$ then folds these
operators backward from the terminal function.

Export the following contracts:
\begin{itemize}[leftmargin=2em]
\item continuity in $(\beta,h,s,\lambda,v,x_1,x_2)$;
\item differentiability once and twice in $\lambda$;
\item differentiation under each finite Gaussian expectation;
\item $|\partial_\lambda U|\le1$ and
      $0\le\partial_{\lambda\lambda}U\le1$;
\item the additive splitting of $U_s^{0,v}$ above the relevant breakpoint;
\item the diagonal backward-Kolmogorov representation used in
      \eqref{eq:GTlinearizedlower}.
\end{itemize}
A finite recursion is important: all regularity can be proved by induction on
the number of semigroup pieces.  No general viscosity-solution theory is
needed for this paper.

\subsection{Specialized GT interpolation}

Formalize Appendix~\ref{app:specialGT} before proving coercivity.  The key
convexity remainder should be defined once:
\begin{lstlisting}
def xi (p : RSPoint) (s u : Real) : Real :=
  p.beta ^ 2 * (1 - s) * p.q * u +
    s * p.beta ^ 2 / 2 * u ^ 2

def xiRemainder (p : RSPoint) (s x y : Real) : Real :=
  xi p s x - xi p s y - deriv (xi p s) y * (x - y)

theorem xiRemainder_eq (p : RSPoint) :
    xiRemainder p s x y = s * p.beta ^ 2 / 2 * (x - y) ^ 2 := by
  simp [xiRemainder, xi]
  ring
\end{lstlisting}
From $s\ge0$, obtain nonnegativity immediately.  This identity supplies the
sign of the GT interpolation derivative.

The auxiliary common Gaussian added in the proof changes covariance but not
the normalized quenched pressure.  State both facts:
\[
 \operatorname{Cov}(H^{\rm cen}+G)
 =N\xi_s(R),
 \qquad
 \frac1N\E\log\sum_\sigma e^{H^{\rm cen}(\sigma)+G}
 =
 \frac1N\E\log\sum_\sigma e^{H^{\rm cen}(\sigma)}.
\]
The second equality uses $\E G=0$ after pulling the common factor $e^G$ out
of the partition sum.

The interpolation theorem should expose every endpoint and derivative
identity:
\begin{lstlisting}
theorem specialized_gt_bound
    (p : RSPoint) (hN : 0 < N)
    (hs : s ∈ Set.Icc (0 : Real) 1)
    (hv : v ∈ overlapRange N) (lambda : Real) :
    disorderE (fun w =>
      Real.log (constrainedPartition2 (pathH p N s w) v)) / N
      <= gtFunctional p s lambda v := by
  have hzero := gt_endpoint_zero ...
  have hone := gt_endpoint_one ...
  have hmono : forall t ∈ Set.Ioo (0 : Real) 1,
      deriv interpolationPressure t <= 0 := by
    ...
  exact hzero.trans (by linarith [monotone_of_deriv_nonpos hmono, hone])
\end{lstlisting}
Do not conceal normalization in the endpoint lemmas.  In particular, verify
the factors $1/N$, $1/2$, the two replicas, and the correction
$\mathcal L_s(v)$ separately.

\subsection{Flatness and the multiplier derivative}

The identity
\[
 \mathfrak P_s(0,v)=2P_s^*,\qquad 0\le v\le1,
\]
comes from the additive semigroup solution.  Formalize it before taking
derivatives.  It prevents unnecessary differentiation of expressions that
are already constant in $v$.

Let
\[
 W_s^v=\left.\partial_\lambda U_s^{\lambda,v}\right|_{\lambda=0}.
\]
The upper-interval product formula
\[
 W_s^v(u,x_1,x_2)=\Psi_x(u,x_1)\Psi_x(u,x_2)
\]
should be proved by checking that both sides satisfy the same finite
semigroup recursion, not by appealing to uniqueness of a large PDE class.
On the diagonal below $v$, the finite-semigroup conditional expectation
gives
\[
 \partial_\lambda\mathfrak P_s(0,v)=g_s(v)-v.
\]
Store this as the bridge theorem between the scalar AT sign and GT
coercivity.

For signed overlaps $-q\le v\le q$, define jointly Gaussian variables with
common variance $\beta^2q$ and covariance
$\beta^2(1-s)q+s\beta^2v$.  Before constructing them, prove the covariance
matrix is positive semidefinite.  Gaussian covariance differentiation gives
\[
 f_s'(v)=s\beta^2\E[
 \operatorname{sech}^2Y_1(v)\operatorname{sech}^2Y_2(v)].
\]
Cauchy--Schwarz and equality of marginal laws yield $f_s'(v)\le s\alpha$.
The proof requires the square-integrability identity
\[
 \E[\operatorname{sech}^4Y_i(v)]=\alpha/\beta^2.
\]

Define the combined signed derivative $\eta_s$ exactly as in
\eqref{eq:GTeta}.  Prove branch agreement at $q$ before proving continuity.
The local coercive sign should be stated in multiplication form
\[
 \eta_s(v)(v-q)\le-c|v-q|^2,
\]
because this is the form consumed by the choice of Lagrange multiplier.

\subsection{Uniform Taylor coercivity}

Lemma~\ref{lem:Taylorcoercivity} is deterministic and should be formalized
early.  A useful Lean statement keeps all constants explicit:
\begin{lstlisting}
theorem uniform_taylor_coercivity
    (P : Type*) [TopologicalSpace P] [CompactSpace P]
    (q : P -> Real) (H : P -> Real -> Real -> Real)
    (lambda0 M c kappa : Real)
    (hlambda0 : 0 < lambda0) (hM : 0 < M)
    (hsecond : forall p lambda v,
      abs lambda <= lambda0 ->
      abs (secondDerivInLambda H p lambda v) <= M)
    (hlocal : forall p v,
      abs (v - q p) <= lambda0 * M / c ->
      dLambda H p 0 v * (v - q p) <= -c * (v - q p) ^ 2)
    (hfar : forall p v,
      lambda0 * M / c <= abs (v - q p) ->
      infOverLambda (H p) v <= H p 0 (q p) - kappa) :
    exists c0 > 0, forall p v,
      infOverLambda (H p) v
        <= H p 0 (q p) - c0 * (v - q p) ^ 2 := by
  ...
\end{lstlisting}
The exact neighborhood parameter may be adjusted to simplify the proof.
The local multiplier is
$\lambda=-\partial_\lambda H_p(0,v)/M$.  Prove that it lies in the allowed
interval before applying Taylor's theorem.  For the far region, use
$|v-q_p|\le2$ to convert the fixed loss into a quadratic loss.

Keep the signed far-gap proof in separate cases:
\begin{itemize}[leftmargin=2em]
\item $v\in[-q,1]$ away from $q$, using continuity and the strict sign;
\item $v<-q$ with $s>0$, using strict Cauchy--Schwarz in
      \eqref{eq:signedhopfcole}--\eqref{eq:signedpositivity};
\item $v\le-q$ with $s$ near zero, using \eqref{eq:GTendpoint};
\item $s$ bounded away from zero, using compactness.
\end{itemize}
For strict Cauchy--Schwarz, prove that equality would force proportionality
almost everywhere and then contradict it on a set of positive Gaussian
measure.  A pointwise inequality with no equality analysis cannot supply the
strict compact gap.

The output of this block is Proposition~\ref{prop:GTcoercivity}, preferably
in two forms:
\[
 \frac1N\E\log Z_{N,s}^{(2)}(v)
 \le2P_s^*-c(v-q)^2
\]
for attainable $v$, and an exponential-weight form ready for summing over
$v\in\mathcal R_N$.

\subsection{Uniform second derivative in the multiplier}

The paper uses a uniform bound on
$\partial_{\lambda\lambda}\mathfrak P_s$.  The finite-spin terminal function
has second derivative at most one.  Under the nonlinear semigroup, the
second derivative is a conditional variance plus the semigroup average of
the terminal second derivative, hence remains bounded by a constant
depending only on the total mass of the finite recursion.

Formalize a generic preservation lemma:
\begin{lstlisting}
theorem log_gaussian_semigroup_second_param_bound
    (h1 : forall x, abs (dParam F x) <= L1)
    (h2 : forall x, abs (d2Param F x) <= L2)
    (hm : m ∈ Set.Icc (0 : Real) mMax) :
    forall x,
      abs (d2Param (logGaussianSemigroup m D F) x)
        <= L2 + mMax * L1 ^ 2 := by
  ...
\end{lstlisting}
Inducting over the finite path gives \eqref{eq:GTsecondderivative}.  This
approach records the source of the constant and avoids a bare compactness
argument over an unbounded $\lambda$ parameter.

\subsection{Coupled pressure and preliminary concentration}

Define the coupled pressure by finite sums and prove the elementary
log-sum-exp bound directly from GT coercivity.  The relevant deterministic
lemma is
\[
 \log\sum_{v\in I}e^{L_v}
 \le \log|I|+\max_{v\in I}L_v.
\]
Since a finite set may be empty in general, either assume nonemptiness or
use the known value $1\in\mathcal R_N$ for $N>0$.

For the Gaussian maximum estimate, first prove that each constrained log
partition $L_v$ is Lipschitz in the full disorder vector with one common
constant of order $\sqrt N$.  Compute the squared norm of its gradient,
including both the pair-coupling coordinates and the site-field coordinates.
The two replicas introduce a factor that must be retained.  Apply the finite
maximum concentration theorem with $|\mathcal R_N|\le N+1$ to obtain
\eqref{eq:gaussianmax}.

The smart-path differentiation identity should be a standalone theorem:
\begin{lstlisting}
theorem freeEnergyGap_deriv
    (p : RSPoint) (hN : 0 < N)
    (hs : s ∈ Set.Ioo (0 : Real) 1) :
    HasDerivAt (freeEnergyGap p N) (p.beta ^ 2 / 4 * momentA p N s) s := by
  -- Differentiate the finite Gibbs sum, then apply Gaussian IBP.
  ...
\end{lstlisting}
A companion theorem extends the equality continuously to the closed
interval.  Check the diagonal $i=j$ contribution separately; it is a common
source of an erroneous $O(N^{-1})$ term.

From convexity of the exponential,
\[
 \frac{\lambda_0}{2}A_s
 \le F_{N,s}(\lambda_0).
\]
Combine this with the pressure envelope and
$D_N'(s)=\beta^2A_s/4$ to obtain \eqref{eq:gronwallDN}.  Use an interval
Gronwall theorem with initial condition $D_N(0)=0$.  The output should keep
the explicit rate
\[
 \epsilon_N=C\left(\sqrt{\frac{\log(N+1)}N}
              +\frac{\log(N+1)}N\right),
\]
and prove separately that $\epsilon_N\to0$.

\subsection{Fixed-deviation tails}

Let
\[
 Y_{N,s}=\log\langle
   e^{\lambda_0NQ_{12}^2/2}\rangle_s.
\]
Prove that it is nonnegative and that its expectation is
$2NF_{N,s}(\lambda_0)$.  Its disorder Lipschitz constant is again of order
$\sqrt N$.  Gaussian concentration then gives an exceptional event with
probability exponentially small in $N$.

Formalize the elementary implication outside that event:
\[
 |Q_{12}|\ge\varepsilon
 \quad\Longrightarrow\quad
 \1_{\{|Q_{12}|\ge\varepsilon\}}
 \le
 \exp\left(-\frac{\lambda_0\varepsilon^2N}{2}
     +Y_{N,s}\right).
\]
This is a pointwise statement on the finite replica space.  Taking Gibbs and
disorder expectations comes afterward.  For the finitely many small $N$,
enlarge the multiplicative constant by an explicit finite maximum; do not
use an informal phrase about changing constants without constructing one.

A reusable Lean conclusion is
\begin{lstlisting}
theorem fixed_deviation_tail
    (U : UniformATData) (eps : Real) (heps : 0 < eps) :
    exists c > 0, exists C > 0, forall N, 0 < N ->
      forall x ∈ U.K, forall s ∈ Set.Icc (0 : Real) 1,
      nuAt x N s (fun sigmas =>
        indicator (abs (Q12At x N sigmas) >= eps)) <=
          C * Real.exp (-c * N) := by
  ...
\end{lstlisting}
Use either $>$ or $\ge$ consistently.  The paper switches harmlessly between
them only after constants are adjusted; Lean requires an explicit inclusion.

\subsection{Cavity configuration and endpoint factorization}

For a configuration of size $N$, define the cavity spin
$\varepsilon=\sigma(\operatorname{Fin.last}(N-1))$ and the restriction to
the first $N-1$ coordinates.  A clean implementation uses the equivalence
\[
 \operatorname{Fin}N\simeq \operatorname{Fin}(N-1)\,\sqcup\,\mathbf 1
\]
under $0<N$.  Prove that summing over full configurations is equal to summing
over a restricted configuration and the last Boolean spin.

The centered cavity overlap keeps denominator $N$:
\[
 Q^-_{ab}=\frac1N\sum_{i<N}\sigma_i^a\sigma_i^b-q.
\]
Do not replace this denominator by $N-1$.  Prove the exact relation
\[
 Q_{ab}=Q^-_{ab}+N^{-1}\varepsilon_a\varepsilon_b.
\]
All three diagonal vectors in \eqref{eq:theta} come from this normalization.

At $u=0$, conditional on
$Y=h+\beta\sqrt q\,Z$, cavity spins are independent with mean $\tanh Y$.
Formalize the conditional finite-spin distribution first:
\[
 \Pr(\varepsilon=e\mid Y)
 =\frac{e^{eY}}{2\cosh Y}.
\]
Then prove for any finite set of distinct replica indices $A$,
\[
 \E_0\prod_{a\in A}\varepsilon_a
 =
 \begin{cases}
 1,&|A|=0,\\
 q,&|A|=2,\\
 r,&|A|=4,
 \end{cases}
\]
for the cases needed here.  The more useful general rule reduces a spin
monomial by parity of multiplicities before counting the number of odd
indices.

\subsection{The replica derivative operator}

Represent an overlap edge as an unordered pair of replica indices.  A
project structure can enforce distinct endpoints:
\begin{lstlisting}
structure ReplicaEdge (n : Nat) where
  left : Fin n
  right : Fin n
  ne : Not (left = right)
\end{lstlisting}
Alternatively use ordered pairs and prove symmetry normalization.  The first
choice simplifies parity counting; the second aligns more directly with
matrix indices.  Do not quotient by edge symmetry until the simp behavior
has been tested.

For an observable of $n$ replicas, define the three parts of
$\mathcal D_n$ with explicit embeddings into $n+2$ replicas:
\begin{align*}
\mathcal D_nF={}&
 \sum_{1\le a<b\le n}F\varepsilon_a\varepsilon_bQ^-_{ab}\\
&-n\sum_{a=1}^n
 F\varepsilon_a\varepsilon_{n+1}Q^-_{a,n+1}\\
&+\frac{n(n+1)}2
 F\varepsilon_{n+1}\varepsilon_{n+2}Q^-_{n+1,n+2}.
\end{align*}
The coefficient in the last line counts the normalization replicas produced
by differentiating a Gibbs expectation.  It is not the cardinality of a
displayed edge set, so record it directly.

State the cavity derivative theorem as a \texttt{HasDerivAt} result on
$u\in(0,1)$:
\begin{lstlisting}
theorem cavity_expect_hasDerivAt
    (F : Replicas N n -> Real) (hF : Bounded F)
    (hu : u ∈ Set.Ioo (0 : Real) 1) :
    HasDerivAt
      (fun u => cavityNu p N s u F)
      (s * p.beta ^ 2 * cavityNu p N s u (cavityOp n F)) u := by
  ...
\end{lstlisting}
The observable on the right formally lives on $n+2$ replicas.  Reflect that
in the type of \texttt{cavityOp}; relying on an implicit coercion will make
the second derivative hard to state.

Prove endpoint continuity separately.  The square-root coefficients in the
Hamiltonian are not differentiable at the endpoints, but the covariance
derivative and the final expectation have one-sided continuous extensions.

\subsection{Cavity replacement estimates}

For
\[
 M_3(u)=\nu_{s,u}[|Q^-_{12}|^3],
\]
differentiate with the cavity operator and bound every extra overlap by $2$.
Formalize replica H\"older in a symmetric form:
\[
 \nu[|Q^-_{ab}Q^-_{cd}Q^-_{ef}|]
 \le\prod_{\text{three edges}}
      \nu[|Q^-_{\cdot\cdot}|^3]^{1/3}
 =M_3(u).
\]
The final equality uses replica exchangeability, so all six indices may have
different coincidence patterns.  Verify that the one-edge marginal is
always the same; no independence is claimed.

Apply a two-sided or backward Gronwall lemma to
$|M_3'(u)|\le CM_3(u)$ to get
\eqref{eq:cavityM3comparison}.  Next prove
\[
 |Q^-_{12}|\le |Q_{12}|+N^{-1}
\]
and the real inequality $(x+y)^3\le4(x^3+y^3)$ for nonnegative $x,y$.
This yields \eqref{eq:cavityM3endpoint}.

For products of two overlaps, the derivative adds one overlap.  Integrating
the derivative and applying the triple H\"older bound proves
\eqref{eq:cavityquadraticcomparison}.  Keep the following real Young
inequality as a named lemma:
\begin{lstlisting}
theorem abs_div_nat_le_cube_add (x : Real) (hN : 0 < N) :
    abs x / N <= abs x ^ 3 / 3 + 2 / (3 * (N : Real) ^ (3 / 2 : Real)) := by
  -- Use Young's inequality with exponents 3 and 3/2.
  ...
\end{lstlisting}
Because real powers create extra positivity obligations, an equivalent
version with $N\sqrt N$ in the denominator may be easier:
\[
 \frac{|x|}{N}\le\frac{|x|^3}{3}+\frac{2}{3N^{3/2}}.
\]
Prove the casted $N>0$ facts once and reuse them.

\subsection{Cavity Taylor remainder}

For each centered cavity observable $F$, establish
\[
 \nu_{s,0}[F]=0.
\]
This follows from endpoint factorization and the definitions of $q$ and $r$.
Use the integral Taylor formula
\[
 f(1)=f(0)+f'(0)+\int_0^1(1-u)f''(u)\,\dd u.
\]
A local theorem specialized to $[0,1]$ may be much easier to apply than a
fully general second-order Taylor theorem.

When differentiating twice, the replica count changes after the first
application of $\mathcal D$.  Each summand must carry the number of replicas
on which it truly depends before the operator is applied again.  Define a
small syntax of cavity monomials if direct dependent types become hard:
\begin{lstlisting}
structure CavityMonomial where
  replicas : Nat
  spinEdges : List (ReplicaEdge replicas)
  overlapEdges : List (ReplicaEdge replicas)
  coefficient : Real
\end{lstlisting}
Give this syntax an evaluation function.  Then prove that one derivative
maps each monomial to a finite sum of monomials and adds exactly one overlap
edge.  The second derivative bound reduces to: every produced term has three
overlap edges, bounded spin factor, and a uniformly bounded coefficient.

The diagonal terms already carry $N^{-1}$.  One derivative adds a single
overlap, and the Young estimate converts
$N^{-1}\nu|Q^-|$ into the allowed cubic remainder plus $N^{-3/2}$.

The final cavity system theorem should not use big-$O$ notation:
\begin{lstlisting}
theorem cavity_system
    (U : UniformATData) :
    exists C > 0, forall N, 0 < N ->
      forall x ∈ U.K, forall s ∈ Set.Icc (0 : Real) 1,
      exists rho : Fin 3 -> Real,
        (identity3 - s • cavityL x) *v momentVec x N s =
          (1 / N : Real) • theta x + rho
        ∧ normInf rho <=
          C * ((N : Real)⁻¹ / Real.sqrt N +
            thirdMoment x N s) := by
  ...
\end{lstlisting}
Choose one expression for $N^{-3/2}$ and prove its equivalence to
$1/(N\sqrt N)$ under $N>0$.

\subsection{Finite coefficient verification}

The coefficient table \eqref{eq:cavitytable} is suitable for a fully
decidable check once $q$ and $r$ are treated as formal variables.  Encode a
spin monomial by the parity set of its replica indices.  At the decoupled
endpoint, assign the symbolic value
\[
 \operatorname{mom}(A)=
 \begin{cases}
 1,&|A|=0,\\
 q,&|A|=2,\\
 r,&|A|=4.
 \end{cases}
\]
Multiplication of spin monomials becomes symmetric difference of parity sets.
The factor
\[
 \E_0[(\varepsilon_p\varepsilon_q-q)\varepsilon_e]
\]
then normalizes to $1-q^2$, $q-q^2$, or $r-q^2$ according to whether the
edge is equal, adjacent, or disjoint.

There should be two checks of the table:
\begin{itemize}[leftmargin=2em]
\item a general parity lemma proving the three symbolic values;
\item a finite \texttt{native\_decide} or \texttt{decide} computation for
      the edge multiplicities in the $n=2,3,4$ cases.
\end{itemize}
The computational check counts combinatorial coefficients only.  The
mathematical theorem connects those counts to cavity expectations.  Keeping
both makes the calculation easy to audit.

Expected rows are
\[
 (b_2,-4b_1,3b_0),
\]
\[
 (b_1,b_2-2b_1-3b_0,6b_0-3b_1),
\]
and
\[
 (b_0,4b_1-8b_0,b_2-8b_1+10b_0).
\]
After the finite count, use \texttt{ring} to collect the symbolic $b_i$
terms.

\subsection{The three-dimensional cavity matrix}

Use vectors \texttt{Fin 3 -> Real} and matrices
\texttt{Matrix (Fin 3) (Fin 3) Real}.  Define the entries in one location:
\begin{lstlisting}
def cavityB2 (p : RSPoint) : Real :=
  p.beta ^ 2 * (1 - p.q ^ 2)

def cavityB1 (p : RSPoint) : Real :=
  p.beta ^ 2 * p.q * (1 - p.q)

def cavityB0 (p : RSPoint) : Real :=
  p.beta ^ 2 * (p.r - p.q ^ 2)

def theta (p : RSPoint) : Fin 3 -> Real :=
  ![1 - p.q ^ 2, p.q - p.q ^ 2, p.r - p.q ^ 2]

def cavityL (p : RSPoint) : Matrix (Fin 3) (Fin 3) Real :=
  ![![cavityB2 p, -4 * cavityB1 p, 3 * cavityB0 p],
    ![cavityB1 p,
       cavityB2 p - 2 * cavityB1 p - 3 * cavityB0 p,
       6 * cavityB0 p - 3 * cavityB1 p],
    ![cavityB0 p,
       4 * cavityB1 p - 8 * cavityB0 p,
       cavityB2 p - 8 * cavityB1 p + 10 * cavityB0 p]]
\end{lstlisting}
The \texttt{![...]} notation may need the relevant matrix and vector
imports.  Keep entrywise simp lemmas so \texttt{fin\_cases i} followed by
\texttt{ring} can prove matrix equalities.

Define $S$ and an explicit $S^{-1}$.  Prove both
$SS^{-1}=I$ and $S^{-1}S=I$ by extensionality and finite cases.  Then prove
\eqref{eq:triangular} entrywise.  Do not rely only on a determinant
calculation: the upper-right Jordan entry matters for the uniform inverse
bound.

The scalar inequalities required for stability are
\[
r\le q,\qquad
\alpha-\beta^2\gamma'=2\beta^2(q-r)\ge0,
\]
and hence
\[
 \beta^2\gamma'\le\alpha<1.
\]
For $s\in[0,1]$, prove
\[
1-s\alpha\ge1-\alpha>0
\]
when $\alpha\ge0$; nonnegativity of $\alpha$ follows from the sech-fourth
formula.  For $\gamma'$, split on its sign exactly as in the paper.

One route to the inverse bound is to write the inverse of the triangular
matrix explicitly:
\[
 \begin{pmatrix}
 a&b&0\\0&a&0\\0&0&c
 \end{pmatrix}^{-1}
 =
 \begin{pmatrix}
 a^{-1}&-b/a^2&0\\0&a^{-1}&0\\0&0&c^{-1}
 \end{pmatrix}.
\]
Here
\[
a=1-s\beta^2\gamma',\qquad
b=-s\beta^2\gamma'',\qquad
c=1-s\alpha.
\]
Compactness bounds $\beta$, $q$, and $r$, so it also bounds $b$.  Conjugating
by the fixed matrices $S$ and $S^{-1}$ supplies a concrete infinity-norm
bound.  This avoids relying on continuity of matrix inversion, though that
is a valid alternative after invertibility is proved.

The replicon row can be checked without inverting the full matrix:
\begin{lstlisting}
def repliconRow : Fin 3 -> Real := ![1, -2, 1]

theorem repliconRow_mul_L (p : RSPoint) :
    Matrix.vecMul repliconRow (cavityL p) =
      p.alpha • repliconRow := by
  funext j
  fin_cases j <;> simp [repliconRow, cavityL, RSPoint.alpha,
    cavityB0, cavityB1, cavityB2] <;> ring

theorem repliconRow_dot_theta (p : RSPoint) :
    dotProduct repliconRow (theta p) = p.alpha / p.beta ^ 2 := by
  have hb : p.beta ^ 2 != 0 := pow_ne_zero 2 (ne_of_gt p.beta_pos)
  simp [repliconRow, theta, RSPoint.alpha]
  field_simp
  ring
\end{lstlisting}
Depending on simplification, it may be better to prove the equivalent
identity
\[
 \beta^2\bigl(\texttt{repliconRow dot theta}\bigr)=\alpha
\]
first, then divide using $\beta>0$.

\subsection{Deterministic cubic absorption}

Abstract the final closure from all probability definitions.  The needed
lemma has the following shape.  Suppose $A_N\ge0$,
\[
 A_N\le aN^{-1}+bN^{-3/2}+dT_N,
\]
and for every fixed $\varepsilon>0$,
\[
 T_N\le\varepsilon A_N+C_\varepsilon e^{-c_\varepsilon N}.
\]
Choose $\varepsilon>0$ with $d\varepsilon\le1/2$ and move the resulting term
to the left.  Then bound
\[
 N^{-3/2}\le N^{-1},\qquad
 e^{-cN}\le C_cN^{-1}
\]
for $N\ge1$.  The latter can use boundedness of
$x\mapsto xe^{-cx}$ on $[0,\infty)$.

The split behind the second assumption is pointwise:
\[
 |Q|^3
 \le\varepsilon Q^2+
 8\1_{\{|Q|>\varepsilon\}}
 \quad\text{when }|Q|\le2.
\]
Prove it by cases on $|Q|\le\varepsilon$.  In the first case,
$|Q|^3\le\varepsilon Q^2$; in the other case, use $|Q|^3\le8$.
This small lemma should not depend on probability.

A suggested interface is
\begin{lstlisting}
theorem absorb_cubic_remainder
    (A T : Nat -> Real) (a b d : Real)
    (hA : forall N, 0 < N -> 0 <= A N)
    (hpre : forall N, 0 < N ->
      A N <= a / N + b / ((N : Real) * Real.sqrt N) + d * T N)
    (htail : forall eps, 0 < eps ->
      exists C c, 0 < C ∧ 0 < c ∧ forall N, 0 < N ->
        T N <= eps * A N + C * Real.exp (-c * N)) :
    exists C, 0 < C ∧ forall N, 0 < N -> A N <= C / N := by
  ...
\end{lstlisting}
In the uniform application, all functions also depend on the compact
parameter and on $s$; either include those indices in the abstract lemma or
apply it to suprema.

\subsection{Free-energy conclusion}

Once
\[
 A_s\le C_{\mathcal K}/N
\]
is available, the free-energy result is short.  Integrate the exact identity
\[
 D_N'(s)=\frac{\beta^2}{4}A_s
\]
from zero to one.  Prove $A_s\ge0$ from positivity of Gibbs expectation, so
$D_N(1)\ge0$.  For the upper bound, use a uniform upper bound on $\beta^2$
over $\mathcal K$ and the fact that the interval has length one.

The endpoint equality
\[
P_1^*=\phi^{\rm RS}(\beta,h)
\]
and $\phi_N(1)=\phi_N(\beta,h)$ should be explicit rewrite lemmas.  With
them, the final proof should require no unfolding of Hamiltonians.

\subsection{Replicon susceptibility}

Let $D_s=A_s-2B_s+C_s$.  Apply the row vector $(1,-2,1)$ to the cavity
system.  The two finite algebra lemmas above give
\[
 (1-s\alpha)D_s
 =\frac{\alpha}{\beta^2N}+\rho^{(R)}_{N,s}.
\]
Bound
\[
 |\rho^{(R)}_{N,s}|
 \le \|\,(1,-2,1)\,\|_1\|\rho_{N,s}\|_\infty
 =4\|\rho_{N,s}\|_\infty.
\]
Record the factor four rather than absorbing it immediately.

To show the scaled cubic moment vanishes uniformly, use
\[
 N\nu_s[|Q|^3]
 \le C\varepsilon+
 8NC_{\varepsilon}e^{-c_{\varepsilon}N}.
\]
For each fixed $\varepsilon$, the exponential term tends to zero uniformly
in the parameters.  Then send $\varepsilon$ to zero.  Formalize this as an
$\varepsilon$ proof of uniform convergence, or define
\[
 e_N=\sup_{(\beta,h,s)\in\mathcal K\times[0,1]}
      N\nu_s[|Q|^3]
\]
and prove $e_N\to0$ using a limsup inequality.

Finally divide by $1-s\alpha$.  The common AT gap makes the denominator
uniformly positive.  State the result as uniform convergence:
\begin{lstlisting}
theorem replicon_susceptibility_uniform
    (U : UniformATData) :
    TendstoUniformlyOn
      (fun N x_s =>
        N * (momentAAt x_s N - 2 * momentBAt x_s N +
          momentCAt x_s N))
      (fun x_s =>
        alphaAt x_s.1 /
          (betaAt x_s.1 ^ 2 * (1 - x_s.2 * alphaAt x_s.1)))
      atTop (U.K ×ˢ Set.Icc (0 : Real) 1) := by
  ...
\end{lstlisting}
If Mathlib's uniform-convergence predicate is inconvenient, use the direct
quantifier form from Theorem~\ref{thm:main}:
for every $\varepsilon>0$, there is $N_0$ such that the absolute difference
is at most $\varepsilon$ for every $N\ge N_0$, every parameter in
$\mathcal K$, and every $s\in[0,1]$.

\subsection{Uniformity on compact parameter sets}

Do not use a symbol $C_{\mathcal K}$ without recording what it depends on.
A robust pattern is to prove lemmas of the form
\begin{lstlisting}
exists C > 0, forall x ∈ K, forall s ∈ Set.Icc (0 : Real) 1, ...
\end{lstlisting}
and assign descriptive names to the witnesses.  Recommended names are
\begin{itemize}[leftmargin=2em]
\item \texttt{betaBound}, \texttt{qLower}, and \texttt{atGap};
\item \texttt{gtSecondBound} and \texttt{matrixInverseBound}.
\end{itemize}

The compact facts needed from the parameter map are:
\begin{itemize}[leftmargin=2em]
\item continuity of $q(\beta,h)$;
\item continuity of $r$ and $\alpha$ by dominated convergence;
\item a positive minimum of $\beta$ and $h$;
\item a positive minimum of $q$;
\item a maximum of $q$ strictly below one;
\item a positive minimum of $1-\alpha$;
\item finite maxima of $\beta$, $q$, $r$, and all continuous semigroup
      derivative bounds.
\end{itemize}
For a strictly positive continuous function on a compact nonempty set,
write a helper theorem returning a positive lower bound.  State
nonemptiness explicitly.  If $\mathcal K$ is empty, the main uniform theorem
is vacuous and may be handled in a separate case.

The proof that $\sup_{\mathcal K}q<1$ requires more than pointwise $q<1$:
continuity makes $q$ attain a maximum, and the pointwise strict inequality at
the maximizer makes that maximum less than one.  The same pattern yields a
strict uniform AT gap.

\subsection{Paper-to-Lean theorem ledger}

The following ledger should be kept current as declarations are completed.
\begin{center}
\begin{tabular}{p{0.25\textwidth}p{0.31\textwidth}p{0.34\textwidth}}
\toprule
paper item & proposed Lean declaration & principal prerequisites\\
\midrule
\eqref{eq:GIBP} &
\texttt{gaussian\_ibp\_fin} &
finite Gaussian vector, differentiability, integrability\\
\eqref{eq:gaussianmaxgeneral} &
\texttt{gaussian\_max\_bound} &
coordinate Lipschitz bound, finite nonempty index set\\
\eqref{eq:PDEupper}--\eqref{eq:PDElower} &
\texttt{psi\_explicit}, \texttt{psi\_solves} &
Gaussian convolution derivatives\\
Lemma~\ref{lem:diffusioncomparison} &
\texttt{diffusion\_sech4\_comparison} &
It\^o formula, time change, density comparison\\
Proposition~\ref{prop:pathRS} &
\texttt{gs\_sign}, \texttt{gs\_uniform\_linear\_gap} &
scalar PDE, diffusion comparison, compactness\\
Lemma~\ref{lem:GTcontinuity} &
\texttt{gtFunctional\_continuous} &
finite-semigroup parameter continuity\\
Lemma~\ref{lem:Taylorcoercivity} &
\texttt{uniform\_taylor\_coercivity} &
one-variable Taylor theorem, compact fixed gap\\
Lemma~\ref{lem:specialGT} &
\texttt{specialized\_gt\_bound} &
Gaussian covariance interpolation, finite recursion\\
Proposition~\ref{prop:GTcoercivity} &
\texttt{gt\_quadratic\_coercivity} &
GT bound, multiplier derivative, signed gaps\\
Lemma~\ref{lem:sublinearpressure} &
\texttt{coupledPressure\_envelope} &
coercivity, Gaussian maximum\\
\eqref{eq:DNprime} &
\texttt{freeEnergyGap\_deriv} &
finite Gibbs differentiation, Gaussian IBP\\
Corollary~\ref{cor:tail} &
\texttt{fixed\_deviation\_tail} &
coupled pressure, Gaussian concentration\\
\eqref{eq:cavityderivative} &
\texttt{cavity\_expect\_hasDerivAt} &
covariance differentiation, replica normalization\\
Proposition~\ref{prop:cavity} &
\texttt{cavity\_system} &
endpoint moments, Taylor remainder, coefficient count\\
\eqref{eq:triangular} &
\texttt{cavityL\_triangular} &
finite $3\times3$ matrix algebra\\
\eqref{eq:inverse} &
\texttt{cavity\_inverse\_uniform} &
strict AT gap, explicit triangular inverse\\
\eqref{eq:preabsorb}--\eqref{eq:thirdsplit} &
\texttt{overlap\_second\_moment\_bound} &
matrix inverse, fixed tail, cubic absorption\\
\eqref{eq:freeenergyidentity} &
\texttt{freeEnergy\_error\_bound} &
smart-path derivative, second moment\\
\eqref{eq:repliconeq} &
\texttt{replicon\_susceptibility\_uniform} &
replicon row, vanishing scaled cubic moment\\
\bottomrule
\end{tabular}
\end{center}

Every ledger row should eventually have a source file and a theorem name
that Lean accepts.  If one paper statement is split into several
declarations, list the public theorem that reassembles it.

\subsection{Suggested public theorem statements}

A pointwise finite-volume theorem may have the following shape:
\begin{lstlisting}
theorem overlap_second_moment_pointwise
    (p : RSPoint) (hat : p.StrictAT) :
    exists C > 0, forall N, 0 < N ->
      forall s ∈ Set.Icc (0 : Real) 1,
        momentA p N s <= C / N := by
  ...
\end{lstlisting}
This statement alone does not imply one common constant on a compact family.
The uniform theorem should quantify in the following order:
\begin{lstlisting}
theorem overlap_second_moment_uniform
    (K : Set (Real × Real))
    (hKc : IsCompact K)
    (hKsub : K ⊆ strictATRegion) :
    exists C > 0, forall N, 0 < N ->
      forall x ∈ K, forall s ∈ Set.Icc (0 : Real) 1,
        momentAAt x N s <= C / N := by
  ...
\end{lstlisting}
Then state the full paper result as a structure or conjunction containing the
second-moment estimate, free-energy estimate, and uniform susceptibility
limit.  Keeping the first two theorems public is useful for downstream work.

The condition $N\ge1$ is represented by \texttt{0 < N}.  The paper's
$C_{\mathcal K}<\infty$ becomes \texttt{exists C, 0 < C} together with the
bound.  There is no need to encode finiteness separately for a real number.

\subsection{Tactic guidance for fragile calculations}

Use \texttt{ring} or \texttt{ring\_nf} for polynomial identities in
$q,r,\beta,s$.  Use \texttt{field\_simp} only after proving every denominator
nonzero.  For inequalities with squares, \texttt{nlinarith} is effective
once square nonnegativity is in the local context.

For \texttt{Fin 2} and \texttt{Fin 3} matrix goals, the reliable pattern is
\begin{lstlisting}
  ext i j
  fin_cases i <;> fin_cases j
  all_goals simp [cavityL, changeMatrix, cavityB0, cavityB1, cavityB2]
  all_goals ring
\end{lstlisting}
If vector notation expands poorly, prove each coordinate as a named lemma
and combine by \texttt{funext}.

For finite Gibbs sums, use equivalences to reindex rather than expanding all
configurations.  Reserve \texttt{native\_decide} for closed finite
combinatorics, never for identities containing arbitrary real parameters.

For measure-theoretic goals, first prove measurability and integrability as
separate lemmas.  A large \texttt{convert} proof that mixes algebra,
measurability, and almost-everywhere equality is hard to maintain.

\subsection{Audit checklist before claiming completion}

A final audit should verify all of the following:
\begin{itemize}[leftmargin=2em]
\item the project builds with a pinned Lean and Mathlib version;
\item no source contains \texttt{sorry}, \texttt{admit}, or a project axiom;
\item every use of division by $N$, $\beta$, $1-s\alpha$, or a partition
      function has a nearby nonzero proof;
\item every logarithm has a positivity proof for its argument;
\item square-root arguments are nonnegative;
\item every Gaussian integral has measurability and integrability evidence;
\item open-interval derivative identities are extended to endpoints only by
      a proved continuity or dominated-convergence theorem;
\item replica embeddings are explicit after every cavity derivative;
\item the $Q^-$ denominator is $N$, not $N-1$;
\item covariance calculations include diagonal corrections;
\item the half-parameter $\gamma_v=x_q/2$ below $|v|$ is present;
\item constrained partition functions are used only at attainable overlaps;
\item all compact uniform constants are selected outside the quantifiers
      over parameters, $s$, and $N$;
\item the final susceptibility statement uses uniform, not merely pointwise,
      convergence;
\item \texttt{\#print axioms} for each public theorem reports only the
      standard logical principles inherited from Mathlib.
\end{itemize}

Useful repository checks are
\begin{lstlisting}
lake build
rg -n "\bsorry\b|\badmit\b|^axiom " SKKusuoka
lake env lean SKKusuoka/Main.lean
\end{lstlisting}
The regular expression should be run on source code, with comments excluded
or reviewed manually.  Also add small executable or example files that
check the $3\times3$ identities and the cavity coefficient enumeration
independently of the analytic development.

\subsection{High-risk proof obligations}

The parts most likely to require new general-purpose library work are the
positive-field uniqueness of the fixed point, the SDE comparison in
Lemma~\ref{lem:diffusioncomparison}, strictness in the signed
Cauchy--Schwarz gap, Gaussian concentration with the exact finite-volume
Lipschitz constant, and the twice-differentiated cavity expectation with
changing replica count.

For each such block, first write the exact theorem signature and prove every
finite or deterministic sublemma around it.  This prevents a hard analytic
lemma from becoming entangled with notation or normalization.  In
particular:
\begin{itemize}[leftmargin=2em]
\item the diffusion block should accept already-proved identities for the
      law at time $q$ and return only the comparison estimate;
\item the strict signed-gap block should accept the explicit two-variable
      Gaussian semigroup and return a continuous negative gap on a compact
      set;
\item the concentration block should accept a Lipschitz function on a finite
      Gaussian vector and return the exact tail bound used later;
\item the cavity second-derivative block should work for a finite list of
      monomials and return a bound by the third overlap moment.
\end{itemize}
Once these interfaces are stable, their proofs can be improved without
changing the main argument.

\subsection{Minimal useful completion stages}

A useful algebraic release consists of the following files:
\begin{itemize}[leftmargin=2em]
\item \texttt{Spins.lean} and \texttt{Gibbs.lean};
\item \texttt{Cavity/Coefficients.lean} and \texttt{Matrix3.lean};
\item \texttt{Absorption.lean}.
\end{itemize}
It checks the exact finite-size conventions and the stability calculation.

A useful analytic-core release adds \texttt{Parameters.lean},
\texttt{ScalarPDE.lean}, and \texttt{Diffusion.lean}.  It should end with a
pointwise strict-AT sign theorem.

A coercivity release adds the finite GT recursion, its interpolation bound,
and the signed compact-gap cases.  Its public endpoint is
Proposition~\ref{prop:GTcoercivity}.

The final release adds preliminary concentration, cavity remainders,
compact uniformity, and all three conclusions of Theorem~\ref{thm:main}.
Each release should build without unfinished declarations; work in progress
can live on a separate branch or in files not imported by the release root.

This organization turns each displayed equation in the paper into a
checkable contract and keeps the eventual Lean code close enough to the
mathematical notation for line-by-line comparison.


\section*{Acknowledgements}
S.K.\ acknowledges support from JSPS KAKENHI Grant Numbers
JP22H00099 and JP23K20801.
S.N.\ acknowledges support from JSPS KAKENHI Grant Numbers
JP24K16937 and JP25K00911.

\section*{Statement on the Use of Artificial Intelligence}
Codex was used to assist with revising the presentation and performing
compilation checks.  The authors take full responsibility for the
content of this work.

\begin{thebibliography}{10}

\bibitem{ChenGT}
W.-K.~Chen,
\newblock
\emph{Variational representations for the Parisi functional and the
two-dimensional Guerra--Talagrand bound},
\newblock \emph{Annals of Probability} 45 (2017), 3929--3966.
\href{https://arxiv.org/abs/1501.06635}{arXiv:1501.06635}.

\bibitem{DeyWu}
P.~S.~Dey and Q.~Wu,
\newblock
\emph{Fluctuation results for multi-species
Sherrington--Kirkpatrick model in the replica symmetric regime},
\newblock \emph{Journal of Statistical Physics} 185 (2021),
Paper No.~22.
\href{https://arxiv.org/abs/2012.13381}{arXiv:2012.13381}.

\bibitem{Ledoux}
M.~Ledoux,
\newblock
\emph{The Concentration of Measure Phenomenon},
\newblock Mathematical Surveys and Monographs, Vol.~89,
American Mathematical Society, Providence, RI, 2001.

\bibitem{Lopatto}
P.~Lopatto,
\newblock
\emph{Replica symmetry up to the de Almeida--Thouless line in the
Sherrington--Kirkpatrick model},
\newblock preprint (2026).
\href{https://arxiv.org/abs/2604.11921}{arXiv:2604.11921}.

\bibitem{Panchenko}
D.~Panchenko,
\newblock
\emph{The Sherrington--Kirkpatrick Model},
\newblock Springer Monographs in Mathematics, Springer, New York, 2013.

\bibitem{Talagrand}
M.~Talagrand,
\newblock
\emph{Mean Field Models for Spin Glasses}, Vols.~I--II,
\newblock Springer, 2011.

\end{thebibliography}


\end{document}
